\documentclass[acmsmall,screen,review,anonymous,nonacm]{acmart}
\usepackage[shortlabels]{enumitem}
\usepackage[utf8]{inputenc}
\usepackage{amsmath}
\usepackage{amsthm}
\usepackage{array}
\usepackage{braket}
\usepackage{mathpartir}
\usepackage{framed}
\usepackage{graphicx}
\usepackage{tikz}
\usepackage{float}
\usepackage{placeins}
\usepackage{url}
\usepackage{comment}
\usepackage{hyperref}
%\usepackage{upgreek}
\usetikzlibrary{shapes.geometric,arrows.meta,positioning,calc,fit}
\usepackage{longtable} % For tables spanning multiple pages
\usepackage{multirow}
\usepackage{hhline}
\usepackage{listings}
\usepackage{tcolorbox}
\usepackage{booktabs}
\tcbset{
  fontupper=\fontsize{9pt}{9.6pt}\selectfont,
  fonttitle=\fontsize{9.1pt}{9.7pt}\selectfont\bfseries,
  left=0.35mm,
  right=0.35mm,
  top=0.35mm,
  bottom=0.35mm,
  boxsep=0.35mm,
  before upper={\setlength{\jot}{0.8pt}\setlength{\arraycolsep}{1pt}}
}
\allowdisplaybreaks
\setlength{\hfuzz}{40pt}
\setlength{\vfuzz}{45pt}
\hbadness=10000
\vbadness=10000
\setlength{\emergencystretch}{2em}
\newtheorem{definition}{Definition}[section]
\setlength{\arrayrulewidth}{1.2pt} % Set line width to 1.2pt
\usetikzlibrary{quantikz2}
\hypersetup{colorlinks=true, citecolor=blue, urlcolor=blue, linkcolor=blue}
\setcounter{tocdepth}{2}
\geometry{a4paper}

\newtheoremstyle{defn_style}      % Name
  {}%                             % Space above
  {}%                             % Space below
  {}%                             % Body font
  {}%                             % Indent amount
  {\bfseries}%                    % Theorem head font
  {.}%                            % Punctuation after theorem head
  { }%                            % Space after theorem head, ' ', or \newline
  {\thmname{#1}\thmnumber{#2}\thmnote{ (#3)}} % Theorem head spec (can be left empty, meaning `normal')

\theoremstyle{defn_style}
\newtheorem{defn}{Definition}[section]
\newtheorem{example}{Example}[section]
\newtheorem{pro}{Proposition}[section]

\newtheorem{pros}{Properties}[section]

\newtheorem{lem}{Lemma}[section]

\newtheorem{lemma}[lem]{Lemma}

\newtheorem{thm}{Theorem}[section]
\newcommand{\op}[2]{\ket{#1}\!\bra{#2}}
\newcommand{\opp}[3]{\ket{#1}\!_{#3}\!\bra{#2}}
\newcommand{\sem}[1]{\llbracket #1 \rrbracket}
\newcommand{\asteq}{\mathrel{*}=}
\newcommand{\supp}{\textup{supp}}
\newcommand{\spn}{span}
\newcommand{\tr}{\textup{tr}}
\newcommand{\prj}{\textup{proj}}

\newcommand{\SKIP}{\ensuremath{\mathbf{skip}}}
\newcommand{\IF}{\ensuremath{\mathbf{if}}}
\newcommand{\THEN}{\ensuremath{\mathbf{then}}}
\newcommand{\ELSE}{\ensuremath{\mathbf{else}}}
\newcommand{\FI}{\ensuremath{\mathbf{fi}}}
\newcommand{\WHILE}{\ensuremath{\mathbf{while}}}
\newcommand{\MEAS}{\ensuremath{\mathbf{meas}}}
\newcommand{\DO}{\ensuremath{\mathbf{do}}}
\newcommand{\OD}{\ensuremath{\mathbf{od}}}
\newcommand{\FOR}{\ensuremath{\mathbf{for}}}
\newcommand{\ERR}{\ensuremath{\mathbf{error}}}
\newcommand{\ASSERT}{\ensuremath{\mathbf{assert}}}
\newcommand{\FREE}{\ensuremath{\mathbf{free}}}

\title{Relational Analysis for Quantum Key Distribution}
\author{Junlin}
\date{Oct 2025}

\begin{document}

\maketitle

\section{Formal System Soundness}

\subsection{Program Model}
\subsubsection{Syntax}

\begin{definition}[Variable Names to Address Indexes]
\[
\overset{\textbf{Quantum variables}}{\overbrace{
\bar q \mapsto [1,\ldots,n_q]
}}
\qquad
\overset{\textbf{Probability registers}}{\overbrace{
p \mapsto [1,\ldots,n_p]
}}
\qquad
\overset{\textbf{Classical registers}}{\overbrace{
x \mapsto [1,\ldots,n_x]
}} .
\]

Addresses are defined as
\[
(\mathsf{Quantum}\quad \mathsf{Int})
\mid
(\mathsf{Probability}\quad \mathsf{Int})
\mid
(\mathsf{Classical}\quad \mathsf{Int}) .
\]

Descriptor names include
\[
f \mapsto [1,\ldots,n_f],
\qquad
\iota:\mathsf{Int}.
\]
They are not register addresses.
\end{definition}


\begin{definition}[Logical Lists, Length, and Descriptor Views]
The physical register memories are fixed address-indexed memories. Addresses
are not created, deleted, copied, or moved by descriptor commands.

A descriptor store \(\xi\) assigns to each logical list name \(g\) a descriptor
view
\[
\operatorname{desc}_{\xi}(g)=(\ell_g,\alpha_g),
\]
where
\[
\ell_g\in\mathbb N
\]
and
\[
\alpha_g:\{1,\ldots,\ell_g\}\to \mathsf{addresses}.
\]

A descriptor view is mutable. A descriptor command may overwrite the descriptor
currently stored under a logical list name. Such an overwrite changes only how
future logical references are resolved. It does not change the register payload
stored at any physical address.

A reference \(g[i]\) is well formed at \(\xi\) only when
\[
\operatorname{desc}_{\xi}(g)=(\ell_g,\alpha_g)
\qquad
\text{and}
\qquad
1\le i\le \ell_g.
\]
It resolves to
\[
\mathsf{phys}_{\xi}(g[i])=\alpha_g(i).
\]

For lists,
\[
\mathsf{phys}_{\xi}(g[1,\ldots,n])
=
(\alpha_g(1),\ldots,\alpha_g(n)),
\]
provided
\[
n\le \ell_g.
\]

A branch separates descriptor data from register data:
\[
\text{descriptor data: filters, descriptor integers, and descriptor views,}
\]
\[
\text{register data: payload at physical addresses.}
\]
Classical values, probability values, and quantum states are register data.
\end{definition}


\begin{definition}[Static Logical Indices]
In the core language, every displayed index
\[
i,j,t
\]
in a logical reference such as
\[
g[i],
\qquad
x[i],
\qquad
\bar q[j]
\]
is a fixed positive natural number at the syntax level. It is not a runtime
integer expression and not a descriptor integer.

Thus a core reference
\[
g[i]
\]
is well formed at descriptor store \(\xi\) only when
\[
\operatorname{desc}_{\xi}(g)=(\ell_g,\alpha_g)
\qquad\text{and}\qquad
1\le i\le \ell_g.
\]

Expression-indexed selected-qubit notation such as
\[
\bar q[e_1,\ldots,e_k]\asteq U(e)
\]
is syntax sugar only. It expands into finite case analysis over fixed
natural-number references before denotation.
\end{definition}


\begin{definition}[Syntax]
Statements are divided by command mode.

Descriptor commands update descriptor data:
\[
D ::=
f :=_{\$}\mathbf{shuffle}(m)_N
\mid
\iota := \mathbf{length}(g)
\]
\[
\mid
(g_0,g_1)
:=
\mathbf{select}_{f}^{N}\ \mathbf{from}\ J\ \mathbf{for}\ g
\mid
g
:=
\mathbf{merge}_{f}^{N}\ \mathbf{from}\ J\ \mathbf{for}\ g_0,g_1 .
\]

Descriptor integer expressions are
\[
J ::= N \mid \iota .
\]

Register commands read or update register payloads after resolving logical
references:
\begin{align*}
R ::=~&
\mathbf{skip}
\mid
\mathbf{abort}
\mid
x[i] := E
\mid
z :=_{\$}\mathbf{rand}(n,\varphi)
\mid
z :=_{\$} \textit{rand}_{n}(C_m)
\\
&\mid
\bar q[i] := \ket{0}
\mid
\bar q[i] \asteq U
\mid
x[i] := \mathcal{M}[\bar q[j]]
\mid
R_0;R_1
\\
&\mid
\mathbf{if}\ B \rightarrow R_0,\ \mathbf{else}\rightarrow R_1 .
\end{align*}
Here \(z\) ranges over bit-valued probability register views.

Programs are generated by
\[
S ::=
D
\mid R
\mid S_0;S_1 .
\]

A command
\[
f:=_{\$}\mathbf{shuffle}(m)_N
\]
samples a filter of length \(N\) with exactly \(m\) one-bits.

A command
\[
\iota:=\mathbf{length}(g)
\]
stores the current descriptor length of \(g\).

The commands
\[
(g_0,g_1)
:=
\mathbf{select}_{f}^{N}\ \mathbf{from}\ J\ \mathbf{for}\ g
\]
and
\[
g
:=
\mathbf{merge}_{f}^{N}\ \mathbf{from}\ J\ \mathbf{for}\ g_0,g_1
\]
update descriptor views only.
\end{definition}

\begin{definition}[Classical Expressions and Predicates]
Use classical types
\[
\tau ::= \mathsf{Bit}\mid \mathsf{Int}.
\]

\[
E ::=
a
\mid \iota
\mid x[i]
\mid E_1+E_2 .
\]

\[
B ::=
\top
\mid \bot
\mid E_1=E_2
\mid E_1<E_2
\mid E_1\leq E_2
\mid \neg B
\mid B_1\land B_2
\mid B_1\lor B_2 .
\]

We use \(E_1\neq E_2,\quad E_1>E_2,\quad E_1\geq E_2\)
as standard abbreviations.
\end{definition}


\subsubsection{\textbf{Semantics}}

\begin{definition}[Descriptor Stores and Register States]
A descriptor store \(\xi\) maps filter names to finite bitstrings, descriptor
integer names to integers, and logical list names to descriptor views.

For a logical list \(g\),
\[
\operatorname{desc}_{\xi}(g)=(\ell_g,\alpha_g),
\]
where
\[
\ell_g\in\mathbb N,
\qquad
\alpha_g:\{1,\ldots,\ell_g\}\to \mathsf{addresses}.
\]

A filter descriptor is a finite bit-valued list. Write
\[
\mathsf{BitList}(N)
=
\{0,1\}^{N}
\]
and
\[
\mathsf{Filt}(m,N)
=
\{F\in\{0,1\}^{N}:\operatorname{wt}(F)=m\}.
\]

A c-q state is a finite- or countable-support family
\[
\Delta=(\rho_{\xi})_{\xi},
\]
where each \(\rho_{\xi}\) is a positive sub-normalized physical-register
c-q state and
\[
\sum_{\xi}\operatorname{Tr}(\rho_{\xi})\le 1.
\]

Trace-\(1\) states are admissible initial states. Program denotations may
produce sub-normalized states; in particular, \(\mathbf{abort}\) denotes the
zero state.
\end{definition}


For filter, write
\[
\mathsf{BitList}(N)
=
\{0,1\}^{N}
\]
and
\[
\mathsf{Filt}(m,N)
=
\{F\in\{0,1\}^{N}:\operatorname{wt}(F)=m\}.
\]

A filter expression is either a filter name or a deterministic all-one filter:
\[
\varphi ::= f \mid \mathbf 1_J .
\]
Its interpretation is
\[
\llbracket f\rrbracket_{\xi}=\xi(f),
\]
and
\[
\llbracket \mathbf 1_J\rrbracket_{\xi}
=
1^{\llbracket J\rrbracket_{\xi}}.
\]

The assertion
\[
\Pi_{\varphi}^{N}
\]
means that \(\varphi\) resolves to a bit-valued list of length \(N\). The
assertion
\[
\Pi_{\varphi}^{m,N}
\]
means that \(\varphi\) resolves to a bit-valued list of length \(N\) and
weight \(m\).

\begin{definition}[Resolution and Filter-Indexed Views]
For a logical reference,
\[
\mathsf{phys}_{\xi}(g[i])
=
\alpha_g(i),
\qquad
\text{where }
\operatorname{desc}_{\xi}(g)=(\ell_g,\alpha_g).
\]

For lists,
\[
\mathsf{phys}_{\xi}(g[1,\ldots,n])
=
(\alpha_g(1),\ldots,\alpha_g(n)).
\]

If \(\xi(f)=F\) and
\[
\operatorname{desc}_{\xi}(g)=(\ell_g,\alpha_g),
\]
define
\[
g^{f}
\stackrel{\triangle}{=}
(g|_{F=0},g|_{F=1}),
\]
with relative order preserved inside each filter case. Thus
\[
\mathsf{phys}_{\xi}(g^{f})
=
(\alpha_g(t):1\le t\le \ell_g,\ F_t=0)
\cdot
(\alpha_g(t):1\le t\le \ell_g,\ F_t=1).
\]

For views \(g_0,g_1\), define
\[
(g_0,g_1)^{\uparrow f}
\]
by
\[
\mathsf{phys}_{\xi}((g_0,g_1)^{\uparrow f})
=
\text{the unique interleaving of }
\mathsf{phys}_{\xi}(g_0),
\mathsf{phys}_{\xi}(g_1)
\text{ according to }\xi(f).
\]

Resolution of paired views is componentwise:
\[
\mathsf{phys}_{\xi,\xi'}(u,u')
=
\left(
\mathsf{phys}_{\xi}(u),
\mathsf{phys}_{\xi'}(u')
\right).
\]
\end{definition}

\begin{definition}[Program Semantics]
Descriptor integer expressions are evaluated by
\[
\llbracket N\rrbracket_{\xi}=N,
\qquad
\llbracket \iota\rrbracket_{\xi}=\xi(\iota).
\]

Filter expressions are evaluated by
\[
\llbracket f\rrbracket_{\xi}=\xi(f),
\]
and
\[
\llbracket \mathbf 1_J\rrbracket_{\xi}
=
1^{\llbracket J\rrbracket_{\xi}}.
\]

For
\[
D=f:=_{\$}\mathbf{shuffle}(m)_N,
\]
\[
\sem D(\Delta)_{\xi'}
=
\sum_{\substack{\xi,\ F\in\mathsf{Filt}(m,N)\\
\xi'=\xi[f\mapsto F]}}
\frac{1}{|\mathsf{Filt}(m,N)|}\rho_{\xi}.
\]

For
\[
D=\iota:=\mathbf{length}(g),
\]
if
\[
\operatorname{desc}_{\xi}(g)=(\ell_g,\alpha_g),
\]
then
\[
\sem D(\Delta)_{\xi'}
=
\sum_{\xi:\xi'=\xi[\iota\mapsto \ell_g]}
\rho_{\xi}.
\]

For
\[
D=
(g_0,g_1):=
\mathbf{select}_{f}^{N}\ \mathbf{from}\ J\ \mathbf{for}\ g,
\]
let
\[
e=\llbracket J\rrbracket_{\xi},
\qquad
\xi(f)=F,
\qquad
\operatorname{desc}_{\xi}(g)=(\ell_g,\alpha_g).
\]
The command is well formed only when
\[
e=\ell_g=N,
\qquad
F\in\{0,1\}^{e}.
\]
Then \(D_f(\xi)\) sets
\[
\operatorname{desc}(g_0)
=
\left(
|\{1\le s\le e:F_s=0\}|,
\ t\mapsto (\alpha_g(s):1\le s\le e,\ F_s=0)_t
\right),
\]
and
\[
\operatorname{desc}(g_1)
=
\left(
|\{1\le s\le e:F_s=1\}|,
\ t\mapsto (\alpha_g(s):1\le s\le e,\ F_s=1)_t
\right).
\]
Thus
\[
\sem{D}(\Delta)_{\xi'}
=
\sum_{\xi:D_f(\xi)=\xi'}
\rho_{\xi}.
\]

For
\[
D=
g:=
\mathbf{merge}_{f}^{N}\ \mathbf{from}\ J\ \mathbf{for}\ g_0,g_1,
\]
let
\[
e=\llbracket J\rrbracket_{\xi},
\qquad
\xi(f)=F,
\]
and
\[
\operatorname{desc}_{\xi}(g_0)=(\ell_{g_0},\alpha_{g_0}),
\qquad
\operatorname{desc}_{\xi}(g_1)=(\ell_{g_1},\alpha_{g_1}).
\]
The command is well formed only when
\[
e=N,
\qquad
F\in\{0,1\}^{e},
\]
\[
|\{1\le t\le e:F_t=0\}|=\ell_{g_0},
\qquad
|\{1\le t\le e:F_t=1\}|=\ell_{g_1}.
\]
Then \(D_f(\xi)\) sets \(g\) to the unique descriptor list \(h=(e,\alpha_h)\)
satisfying
\[
\alpha_h(t)=
\begin{cases}
\alpha_{g_0}\!\left(|\{1\le s\le t:F_s=0\}|\right),
&
F_t=0,
\\[0.6ex]
\alpha_{g_1}\!\left(|\{1\le s\le t:F_s=1\}|\right),
&
F_t=1.
\end{cases}
\]
Thus
\[
\sem{D}(\Delta)_{\xi'}
=
\sum_{\xi:D_f(\xi)=\xi'}
\rho_{\xi}.
\]

For
\[
\mathbf{skip},
\]
\[
\sem{\mathbf{skip}}(\Delta)=\Delta.
\]

For
\[
\mathbf{abort},
\]
\[
\sem{\mathbf{abort}}(\Delta)_{\xi}=0
\qquad
\text{for every }\xi.
\]

For sequencing,
\[
\sem{S_0;S_1}
=
\sem{S_1}\circ \sem{S_0}.
\]

For a register command \(R\),
\[
\sem R(\Delta)_{\xi}
=
\mathcal E^R_{\xi}(\rho_{\xi}).
\]

For
\[
z:=_{\$}\mathbf{rand}(n,\varphi),
\]
let
\[
\vec a=\mathsf{phys}_{\xi}(z[1,\ldots,n]).
\]
The register effect is
\[
\mathcal E^R_{\xi}(\rho)
=
2^{-n}
\sum_{s\in\{0,1\}^{n}}
\ket{s}_{\vec a}\bra{s}
\otimes
\operatorname{Tr}_{\vec a}(\rho).
\]

For conditionals, let \(\Pi_{B,\xi}\) be the branchwise classical projector
defined by the predicate \(B\), with descriptor integers interpreted using
\(\xi\). Then
\[
\sem{
\mathbf{if}\ B\rightarrow R_0,\ \mathbf{else}\rightarrow R_1
}(\Delta)_{\xi}
=
\mathcal E^{R_0}_{\xi}
\left(
\Pi_{B,\xi}\rho_{\xi}\Pi_{B,\xi}
\right)
+
\mathcal E^{R_1}_{\xi}
\left(
\Pi_{\neg B,\xi}\rho_{\xi}\Pi_{\neg B,\xi}
\right).
\]
\end{definition}


\begin{definition}[View Index and Local Assertions]
A view index is a displayed left/right view expression
\[
\mathcal V.
\]
It is a proof-side syntactic index.  It contains logical view names, list
expressions, and filter-indexed view expressions such as
\[
g^f,
\qquad
(g_0,g_1)^{\uparrow f}.
\]
It does not contain the current descriptor store.

An assertion is written with an explicit view index:
\[
\Phi_{\mathcal V}.
\]
For every pair of current left/right descriptor stores \((\xi,\xi')\), the
assertion denotes a local projector
\[
\llbracket \Phi_{\mathcal V}\rrbracket_{\xi,\xi'}
\]
on the resolved left/right physical register memories determined by
\[
\mathsf{phys}_{\xi,\xi'}(\mathcal V).
\]
If \(\mathcal V\) is not well formed under \((\xi,\xi')\), then
\[
\llbracket \Phi_{\mathcal V}\rrbracket_{\xi,\xi'}
=
0_{\mathsf{Reg}^{L,R}}.
\]

Thus \(\mathcal V\) is the assertion index, while \(\xi,\xi'\) are semantic
parameters used to interpret that index.

Operators with view subscripts are interpreted locally.  For example,
\[
U_{\bar q[\mathbf i]},
\qquad
\ket a_{x[i]}\bra a,
\qquad
\equiv_{(u,u')}
\]
are interpreted branchwise after resolving the displayed views under
\((\xi,\xi')\).

Descriptor-only assertions are scalar local assertions.  Their denotation is
either
\[
I_{\mathsf{Reg}^{L,R}}
\]
or
\[
0_{\mathsf{Reg}^{L,R}}
\]
under each semantic pair \((\xi,\xi')\).

For filter expressions,
\[
\llbracket \equiv_{(\varphi,\varphi')}\rrbracket_{\xi,\xi'}
\stackrel{\triangle}{=}
\begin{cases}
I_{\mathsf{Reg}^{L,R}},
&
\llbracket\varphi\rrbracket_{\xi}
=
\llbracket\varphi'\rrbracket_{\xi'},
\\[0.4ex]
0_{\mathsf{Reg}^{L,R}},
&
\text{otherwise}.
\end{cases}
\]

For descriptor integers,
\[
\llbracket \equiv_{(\iota,\iota')}\rrbracket_{\xi,\xi'}
\stackrel{\triangle}{=}
\begin{cases}
I_{\mathsf{Reg}^{L,R}},
&
\xi(\iota)=\xi'(\iota'),
\\[0.4ex]
0_{\mathsf{Reg}^{L,R}},
&
\text{otherwise}.
\end{cases}
\]

For one-sided filter length support,
\[
\llbracket \Pi_{\varphi}^{N}\rrbracket_{\xi,\xi'}
\stackrel{\triangle}{=}
\begin{cases}
I_{\mathsf{Reg}^{L,R}},
&
\llbracket\varphi\rrbracket_{\xi}\in\{0,1\}^{N},
\\[0.4ex]
0_{\mathsf{Reg}^{L,R}},
&
\text{otherwise}.
\end{cases}
\]

For one-sided filter weight support,
\[
\llbracket \Pi_{\varphi}^{m,N}\rrbracket_{\xi,\xi'}
\stackrel{\triangle}{=}
\begin{cases}
I_{\mathsf{Reg}^{L,R}},
&
\llbracket\varphi\rrbracket_{\xi}\in\mathsf{Filt}(m,N),
\\[0.4ex]
0_{\mathsf{Reg}^{L,R}},
&
\text{otherwise}.
\end{cases}
\]

For finite-support \(\mu\), define locally
\[
\Pi_z^\mu
\stackrel{\triangle}{=}
\sum_{a\in\operatorname{supp}(\mu)}
\ket a_z\bra a .
\]

For paired \(z,z'\), define locally
\[
\equiv^\mu_{(z,z')}
\stackrel{\triangle}{=}
\sum_{a\in\operatorname{supp}(\mu)}
\ket a_z\bra a
\otimes
\ket a_{z'}\bra a .
\]

For a pair of classical logical references \(z,z'\) of the same finite value
domain \(D\), define locally
\[
\equiv^{D}_{(z,z')}
\stackrel{\triangle}{=}
\sum_{a\in D}
\ket a_z\bra a
\otimes
\ket a_{z'}\bra a .
\]

For probability-register entries \(p[i],p'[j]\),
\[
\equiv^{\mu}_{(p[i],p'[j])}
\stackrel{\triangle}{=}
\equiv^{\operatorname{supp}(\mu)}_{(p[i],p'[j])}.
\]

For a pair of quantum logical references \(\bar q[i],\bar q'[j]\),
\[
\equiv_{(\bar q[i],\bar q'[j])}
\]
is symmetric quantum equality:
\[
=_{\mathrm{sym}}
=
\frac{1}{2}(I\otimes I+\mathsf{SWAP})
\]
on the two resolved systems, extended by identity.

For whole-list equality,
\[
\equiv_{(g,g')}
\]
is interpreted under \((\xi,\xi')\).  If
\[
\operatorname{desc}_{\xi}(g)=(\ell_g,\alpha_g),
\qquad
\operatorname{desc}_{\xi'}(g')=(\ell'_{g'},\alpha'_{g'}),
\qquad
\ell_g=\ell'_{g'},
\]
then
\[
\equiv_{(g,g')}
\stackrel{\triangle}{=}
\bigwedge_{t=1}^{\ell_g}
\equiv_{(g[t],g'[t])}.
\]
If the lengths are not equal, the denotation is \(0\).

For fixed-length list equality,
\[
\equiv_{(g[1,\ldots,n],g'[1,\ldots,n'])}
\]
is interpreted under \((\xi,\xi')\).  If \(n=n'\), then
\[
\equiv_{(g[1,\ldots,n],g'[1,\ldots,n])}
\stackrel{\triangle}{=}
\bigwedge_{t=1}^{n}
\equiv_{(g[t],g'[t])}.
\]
If \(n\neq n'\), the denotation is \(0\).

Filter-indexed view equality is interpreted under \((\xi,\xi')\).  If
\[
\llbracket\varphi\rrbracket_{\xi}=F,
\qquad
\llbracket\varphi'\rrbracket_{\xi'}=F',
\]
then
\[
\equiv_{(g^{\varphi},g'^{\varphi'})}
\]
is nonzero only when
\[
\operatorname{desc}_{\xi}(g)=(\ell_g,\alpha_g),
\qquad
\operatorname{desc}_{\xi'}(g')=(\ell'_{g'},\alpha'_{g'}),
\]
\[
\ell_g=\ell'_{g'},
\qquad
F,F'\in\{0,1\}^{\ell_g},
\qquad
F=F'.
\]
Under these conditions, it is ordinary equality on the resolved ordered views
\[
\mathsf{phys}_{\xi}(g^{\varphi})
\qquad\text{and}\qquad
\mathsf{phys}_{\xi'}(g'^{\varphi'}).
\]
If any condition fails, the denotation is \(0\).

For projective assertions \(\Phi,\Psi\), define \(\Phi\land\Psi\),
\(\Phi\lor\Psi\), and \(\neg\Phi\) branchwise:
\[
\llbracket \Phi\land\Psi\rrbracket_{\xi,\xi'}
=
\operatorname{proj}
\left(
\operatorname{Im}\llbracket \Phi\rrbracket_{\xi,\xi'}
\cap
\operatorname{Im}\llbracket \Psi\rrbracket_{\xi,\xi'}
\right),
\]
\[
\llbracket \Phi\lor\Psi\rrbracket_{\xi,\xi'}
=
\operatorname{proj}
\left(
\operatorname{Im}\llbracket \Phi\rrbracket_{\xi,\xi'}
+
\operatorname{Im}\llbracket \Psi\rrbracket_{\xi,\xi'}
\right),
\]
\[
\llbracket \neg\Phi\rrbracket_{\xi,\xi'}
=
I_{\mathsf{Reg}^{L,R}}
-
\llbracket \Phi\rrbracket_{\xi,\xi'}.
\]
If
\[
\llbracket\Phi\rrbracket_{\xi,\xi'}
\llbracket\Psi\rrbracket_{\xi,\xi'}=0
\]
branchwise, then branchwise
\[
\Phi\lor\Psi=\Phi+\Psi.
\]

The notation
\[
\operatorname{proj}
\left(
\operatorname{Tr}_{\mathcal U}(\Phi)
\right)
\]
means the local register-view projection defined branchwise by
\[
\llbracket
\operatorname{proj}
\left(
\operatorname{Tr}_{\mathcal U}(\Phi)
\right)
\rrbracket_{\xi,\xi'}
=
\operatorname{proj}
\left(
\operatorname{Tr}_{\mathsf{phys}_{\xi,\xi'}(\mathcal U)}
\left(
\llbracket\Phi\rrbracket_{\xi,\xi'}
\right)
\right).
\]
This operation traces only register payload systems.  It does not erase or
modify descriptor data.

The separated conjunction
\[
\Phi_{\mathcal V} * C_{\mathcal U}
\]
is well formed only when the current auxiliary resource contexts certify that
\(\mathcal V\) and \(\mathcal U\) are live and physically disjoint on each
side.  Semantically, it is interpreted branchwise as
\[
\Phi_{\mathcal V} * C_{\mathcal U}
\stackrel{\triangle}{=}
\Phi_{\mathcal V}\land C_{\mathcal U}.
\]

Write \(\equiv\) when the view index is clear.
\end{definition}

\begin{definition}[Approximate Classical-Quantum Lifting]
For c-q states
\[
\Delta=(\rho_{\xi})_{\xi},
\qquad
\Delta'=(\rho'_{\xi'})_{\xi'},
\]
write
\[
\Delta\ \Phi^\delta\ \Delta'
\]
if there exists a witness
\[
\pi=(\omega_{\xi,\xi'})_{\xi,\xi'}
\]
such that
\[
\omega_{\xi,\xi'}
\in
\operatorname{SepD}_{\le 1}(\mathcal H_L:\mathcal H_R),
\]
\[
\sum_{\xi'}
\operatorname{Tr}_R(\omega_{\xi,\xi'})
=
\rho_{\xi},
\]
\[
\sum_{\xi}
\operatorname{Tr}_L(\omega_{\xi,\xi'})
=
\rho'_{\xi'},
\]
and
\[
\sum_{\xi,\xi'}
\operatorname{Tr}
\left(
\llbracket\Phi\rrbracket_{\xi,\xi'}
\omega_{\xi,\xi'}
\right)
\ge
\operatorname{Tr}(\pi)-\delta .
\]

Here
\[
\operatorname{Tr}(\pi)
\stackrel{\triangle}{=}
\sum_{\xi,\xi'}
\operatorname{Tr}(\omega_{\xi,\xi'}).
\]

For \(\delta=0\), write
\[
\Delta\ \Phi\ \Delta' .
\]
\end{definition}

Here
\[
\operatorname{SepD}_{\le 1}(\mathcal H_L:\mathcal H_R)
\]
denotes the set of sub-normalized separable positive operators. That is,
\[
\omega\in\operatorname{SepD}_{\le 1}(\mathcal H_L:\mathcal H_R)
\]
iff
\[
\omega
=
\sum_k p_k\,
\rho^L_k\otimes\rho^R_k,
\qquad
p_k\ge 0,
\qquad
\sum_k p_k\le 1,
\]
for some
\[
\rho^L_k\in D(\mathcal H_L),
\qquad
\rho^R_k\in D(\mathcal H_R).
\]



\subsection{\textbf{Relational Judgment System}}
\begin{definition}[Symbolic Proof Context]
\(\Gamma\) contains persistent symbolic facts: protocol parameters, arithmetic,
and name assumptions. Liveness/freshness/separation are
certified by a resource system, not stored in \(\Gamma\). Entailments have form
\[
\Gamma\vDash \Phi\Longrightarrow\Psi.
\]
\end{definition}

\begin{definition}[Paired Logical Views]
A paired logical view is a finite list
\[
\mathcal V=
[(u_1,u'_1),\ldots,(u_m,u'_m)],
\]
where \(u_i,u'_i\) are corresponding left/right logical references or lists.

Resolution is componentwise:
\[
\mathsf{phys}_{\xi,\xi'}(u,u')
=
\left(
\mathsf{phys}_{\xi}(u),
\mathsf{phys}_{\xi'}(u')
\right).
\]
\end{definition}


\begin{definition}[Classical-Quantum Relational Judgment]
A relational judgment has the form
\[
\Gamma
\vdash
P\sim_{\delta}P'
:
\Phi_{\mathcal V}
\Longrightarrow
\Psi_{\mathcal W}.
\]

It is valid iff for all well-typed \(\Delta,\Delta'\), every witness \(\pi\)
for
\[
\Delta\ \Phi_{\mathcal V}\ \Delta'
\]
with
\[
\operatorname{supp}_{\mathsf{cl}}(\pi)\models\Gamma
\]
satisfies
\[
\llbracket P\rrbracket(\Delta)
\ \Psi_{\mathcal W}^{\delta}\
\llbracket P'\rrbracket(\Delta').
\]

\(\delta\in\mathbb R_{\ge 0}\) bounds postcondition-violating mass. For
\(\delta=0\), write
\[
\Gamma
\vdash
P\sim P'
:
\Phi_{\mathcal V}
\Longrightarrow
\Psi_{\mathcal W}.
\]
\end{definition}


\begin{definition}[Auxiliary-Certified Relational Derivation]
\[
\Gamma\vdash \mathcal U\#\mathcal V
\]
means \(\Gamma\)'s symbolic facts hold and an Auxiliary resource system certifies the
needed left/right liveness and separation.
\end{definition}

\begin{definition}[Command-Mode Footprints]
For every program \(P\), we use four footprint components:
\[
\operatorname{dacc}(P),
\qquad
\operatorname{dupd}(P),
\]
\[
\operatorname{racc}(P),
\qquad
\operatorname{rupd}(P).
\]

\(\operatorname{dacc}(P)\) and \(\operatorname{dupd}(P)\) are the
descriptor-data access and update footprints.

\(\operatorname{racc}(P)\) and \(\operatorname{rupd}(P)\) are the
register-payload access and update footprints.

\[
\operatorname{acc}(P)
\stackrel{\triangle}{=}
\operatorname{dacc}(P)
\cup
\operatorname{racc}(P),
\]
\[
\operatorname{upd}(P)
\stackrel{\triangle}{=}
\operatorname{dupd}(P)
\cup
\operatorname{rupd}(P).
\]

For paired programs,
\[
\operatorname{acc}(P,P')
\stackrel{\triangle}{=}
\operatorname{acc}(P)\cup\operatorname{acc}(P'),
\qquad
\operatorname{upd}(P,P')
\stackrel{\triangle}{=}
\operatorname{upd}(P)\cup\operatorname{upd}(P').
\]

Updates are included in the corresponding access mode:
\[
\operatorname{dupd}(P)\subseteq \operatorname{dacc}(P),
\qquad
\operatorname{rupd}(P)\subseteq \operatorname{racc}(P).
\]

For descriptor integer expressions,
\[
\operatorname{dacc}(N)=\emptyset,
\qquad
\operatorname{dacc}(\iota)=\{\iota\}.
\]

For filter expressions,
\[
\operatorname{dacc}(f)=\{f\},
\qquad
\operatorname{dacc}(\mathbf 1_J)=\operatorname{dacc}(J).
\]

For classical expressions,
\[
\operatorname{dacc}(a)=\emptyset,
\qquad
\operatorname{racc}(a)=\emptyset,
\]
\[
\operatorname{dacc}(\iota)=\{\iota\},
\qquad
\operatorname{racc}(\iota)=\emptyset,
\]
\[
\operatorname{dacc}(x[i])=\{x\},
\qquad
\operatorname{racc}(x[i])=\{x[i]\},
\]
and
\[
\operatorname{dacc}(E_1+E_2)
=
\operatorname{dacc}(E_1)\cup\operatorname{dacc}(E_2),
\]
\[
\operatorname{racc}(E_1+E_2)
=
\operatorname{racc}(E_1)\cup\operatorname{racc}(E_2).
\]

For predicates \(B\), descriptor and register access footprints are defined
recursively from their subexpressions. The constants \(\top,\bot\) have empty
footprints.

For
\[
f:=_{\$}\mathbf{shuffle}(m)_N,
\]
\[
\operatorname{dacc}=\operatorname{dupd}=\{f\},
\qquad
\operatorname{racc}=\operatorname{rupd}=\emptyset.
\]

For
\[
\iota:=\mathbf{length}(g),
\]
\[
\operatorname{dacc}=\{g,\iota\},
\qquad
\operatorname{dupd}=\{\iota\},
\qquad
\operatorname{racc}=\operatorname{rupd}=\emptyset.
\]

For
\[
(g_0,g_1)
:=
\mathbf{select}_{f}^{N}\ \mathbf{from}\ J\ \mathbf{for}\ g,
\]
\[
\operatorname{dacc}
=
\{f,g,g_0,g_1\}
\cup
\operatorname{dacc}(J),
\qquad
\operatorname{dupd}=\{g_0,g_1\},
\qquad
\operatorname{racc}=\operatorname{rupd}=\emptyset.
\]

For
\[
g
:=
\mathbf{merge}_{f}^{N}\ \mathbf{from}\ J\ \mathbf{for}\ g_0,g_1,
\]
\[
\operatorname{dacc}
=
\{f,g,g_0,g_1\}
\cup
\operatorname{dacc}(J),
\qquad
\operatorname{dupd}=\{g\},
\qquad
\operatorname{racc}=\operatorname{rupd}=\emptyset.
\]

For register commands,
\[
\operatorname{dupd}(R)=\emptyset.
\]

For
\[
\mathbf{skip}
\qquad\text{and}\qquad
\mathbf{abort},
\]
all four footprints are empty.

For
\[
x[i]:=E,
\]
\[
\operatorname{dacc}
=
\{x\}\cup\operatorname{dacc}(E),
\]
\[
\operatorname{racc}
=
\{x[i]\}\cup\operatorname{racc}(E),
\qquad
\operatorname{rupd}=\{x[i]\}.
\]

For
\[
z:=_{\$}\mathbf{rand}(n,\varphi),
\]
\[
\operatorname{dacc}
=
\{z\}\cup\operatorname{dacc}(\varphi),
\]
\[
\operatorname{racc}
=
\operatorname{rupd}
=
\{z[1,\ldots,n]\}.
\]

For
\[
z:=_{\$}\textit{rand}_{n}(C_m),
\]
\[
\operatorname{dacc}=\{z\},
\qquad
\operatorname{racc}
=
\operatorname{rupd}
=
\{z[1,\ldots,n]\}.
\]

For
\[
\bar q[i]:=\ket 0,
\]
\[
\operatorname{dacc}=\{\bar q\},
\qquad
\operatorname{racc}=\operatorname{rupd}=\{\bar q[i]\}.
\]

For
\[
\bar q[i]\asteq U,
\]
\[
\operatorname{dacc}=\{\bar q\},
\qquad
\operatorname{racc}=\operatorname{rupd}=\{\bar q[i]\}.
\]

For
\[
x[i]:=\mathcal M[\bar q[j]],
\]
\[
\operatorname{dacc}=\{x,\bar q\},
\]
\[
\operatorname{racc}=\{x[i],\bar q[j]\},
\qquad
\operatorname{rupd}=\{x[i],\bar q[j]\}.
\]

For sequencing,
\[
\operatorname{dacc}(S_0;S_1)
=
\operatorname{dacc}(S_0)\cup\operatorname{dacc}(S_1),
\]
\[
\operatorname{dupd}(S_0;S_1)
=
\operatorname{dupd}(S_0)\cup\operatorname{dupd}(S_1),
\]
\[
\operatorname{racc}(S_0;S_1)
=
\operatorname{racc}(S_0)\cup\operatorname{racc}(S_1),
\]
\[
\operatorname{rupd}(S_0;S_1)
=
\operatorname{rupd}(S_0)\cup\operatorname{rupd}(S_1).
\]

For conditionals,
\[
\operatorname{dacc}
\left(
\mathbf{if}\ B\rightarrow R_0,\ \mathbf{else}\rightarrow R_1
\right)
=
\operatorname{dacc}(B)
\cup
\operatorname{dacc}(R_0)
\cup
\operatorname{dacc}(R_1),
\]
\[
\operatorname{racc}
\left(
\mathbf{if}\ B\rightarrow R_0,\ \mathbf{else}\rightarrow R_1
\right)
=
\operatorname{racc}(B)
\cup
\operatorname{racc}(R_0)
\cup
\operatorname{racc}(R_1),
\]
\[
\operatorname{dupd}
\left(
\mathbf{if}\ B\rightarrow R_0,\ \mathbf{else}\rightarrow R_1
\right)
=
\emptyset,
\]
\[
\operatorname{rupd}
\left(
\mathbf{if}\ B\rightarrow R_0,\ \mathbf{else}\rightarrow R_1
\right)
=
\operatorname{rupd}(R_0)
\cup
\operatorname{rupd}(R_1).
\]
\end{definition}

\begin{definition}[Footprint Equality]
For a mixed footprint set \(A\), the notation
\[
\equiv_A
\]
means the conjunction of the local equalities generated by the typed elements
of \(A\), interpreted under the current view index.

For filter entries,
\[
f\in A
\quad\leadsto\quad
\equiv_{(f,f')}.
\]

For descriptor integer entries,
\[
\iota\in A
\quad\leadsto\quad
\equiv_{(\iota,\iota')}.
\]

For logical descriptor names,
\[
g\in A
\quad\leadsto\quad
\equiv_{(g,g')}.
\]

For register-payload entries,
\[
u\in A
\quad\leadsto\quad
\equiv_{(u,u')},
\]
where \(u,u'\) are resolved using the current view index.

Thus
\[
\equiv_{\operatorname{acc}(P)}
\]
means the local view-index equality needed to run the two copies of \(P\) in
lockstep, and
\[
\equiv_{\operatorname{upd}(P)}
\]
means the corresponding local equality on the updated views.
\end{definition}

\begin{definition}[Descriptor Dependencies of View Indices]
The descriptor dependency set of a view index is written
\[
\operatorname{ddep}(\mathcal V).
\]

For ordinary logical references and lists,
\[
\operatorname{ddep}(g[i])=\{g\},
\qquad
\operatorname{ddep}(g[1,\ldots,n])=\{g\},
\qquad
\operatorname{ddep}(g)=\{g\}.
\]

For descriptor integer expressions,
\[
\operatorname{ddep}(N)=\emptyset,
\qquad
\operatorname{ddep}(\iota)=\{\iota\}.
\]

For filter expressions,
\[
\operatorname{ddep}(f)=\{f\},
\qquad
\operatorname{ddep}(\mathbf 1_J)=\operatorname{ddep}(J).
\]

For filter-induced views,
\[
\operatorname{ddep}(g^{f})
=
\{g,f\}.
\]

For merge-induced views,
\[
\operatorname{ddep}((g_0,g_1)^{\uparrow f})
=
\{g_0,g_1,f\}.
\]

For paired and packaged views, dependencies are collected componentwise.

For a local assertion \(\Phi_{\mathcal V}\),
\[
\operatorname{ddep}(\Phi_{\mathcal V})
\stackrel{\triangle}{=}
\operatorname{ddep}(\mathcal V).
\]
\end{definition}

\textbf{Structural proof rules for QKD:}

\[
\text{\rm (Conseq)}
\quad
\frac{
\Gamma
\vDash
\Phi'_{\mathcal V'}
\Longrightarrow
\Phi_{\mathcal V}
\qquad
\Gamma
\vDash
\Psi_{\mathcal W}
\Longrightarrow
\Psi'_{\mathcal W'}
\qquad
\delta\leq\delta'
\qquad
\Gamma
\vdash
P\sim_{\delta}P':
\Phi_{\mathcal V}
\Longrightarrow
\Psi_{\mathcal W}
}{
\Gamma
\vdash
P\sim_{\delta'}P':
\Phi'_{\mathcal V'}
\Longrightarrow
\Psi'_{\mathcal W'}
}.
\]

\[
\text{\rm (Weaken-\(\Gamma\))}
\quad
\frac{
\Gamma
\vdash
P\sim_{\delta}P':
\Phi\Longrightarrow\Psi
\qquad
\Gamma'\vDash\Gamma
}{
\Gamma'
\vdash
P\sim_{\delta}P':
\Phi\Longrightarrow\Psi
}.
\]

\[
\text{\rm (Seq)}
\quad
\frac{
\Gamma
\vdash
P_0\sim_{\delta}P_1:
\Phi\Longrightarrow\Theta
\qquad
\Gamma
\vdash
P'_0\sim_{\delta'}P'_1:
\Theta\Longrightarrow\Psi
}{
\Gamma
\vdash
P_0;P'_0
\sim_{\delta+\delta'}
P_1;P'_1
:
\Phi\Longrightarrow\Psi
}.
\]

\[
\text{\rm (Skip)}
\quad
\Gamma
\vdash
\mathbf{skip}\sim\mathbf{skip}:
\Phi\Longrightarrow\Phi .
\]

\[
\text{\rm (Abort)}
\quad
\Gamma
\vdash
\mathbf{abort}\sim\mathbf{abort}:
\Phi\Longrightarrow\Psi .
\]

\[
\text{\rm (Ident)}
\quad
{
\Gamma
\vdash
P\sim P:
\equiv_{\operatorname{acc} (P)}
\Longrightarrow
\equiv_{\operatorname{upd} (P)}
}.
\]

\[
\text{\rm (Comp)}
\quad
\frac{
\Gamma
\vdash
P\sim_{\delta_1}P':
\equiv_{\mathcal V}
\Longrightarrow
\equiv_{\mathcal V}
\qquad
\Gamma
\vdash
P'\sim_{\delta_2}P'':
\equiv_{\mathcal V}
\Longrightarrow
\equiv_{\mathcal V}
}{
\Gamma
\vdash
P\sim_{\delta_1+\delta_2}P'':
\equiv_{\mathcal V}
\Longrightarrow
\equiv_{\mathcal V}
}.
\]

An assertion \(C_{\mathcal U}\) is indexed by \(\mathcal U\). 

\[
\text{\rm (Classical-Frame${}^{*}$)}
\quad
\frac{
\Gamma
\vdash
P\sim_{\delta}P':
\Phi_{\mathcal V}
\Longrightarrow
\Psi_{\mathcal W}\qquad
\Gamma\vdash
\mathcal U
\#
\operatorname{upd}(P,P')}{
\Gamma
\vdash
P\sim_{\delta}P':
\Phi_{\mathcal V} * C_{\mathcal U}
\Longrightarrow
\Psi_{\mathcal W} * C_{\mathcal U}
}.
\]

\[
\text{\rm (Quantum-Frame${}^{*}$)}
\quad
\frac{
\begin{array}{c}
\Gamma
\vdash
P\sim_{\delta}P':
\Phi_{\mathcal V}
\Longrightarrow
\Psi_{\mathcal W}
\qquad
\Gamma\vdash
\mathcal U
\#
\operatorname{acc}(P,P')
\end{array}
}{
\Gamma
\vdash
P\sim_{\delta}P':
\Phi_{\mathcal V} * C_{\mathcal U}
\Longrightarrow
\Psi_{\mathcal W} * C_{\mathcal U}
}.
\]


\[
\text{\rm (Swap)}
\quad
\frac{
\Gamma\vdash
\operatorname{acc}(P_0)\#
\operatorname{acc}(P_1)
}{
\Gamma
\vdash
P_0;P_1
\sim
P_1;P_0
:
\equiv_{\mathcal V}
\Longrightarrow
\equiv_{\mathcal V}
}.
\]

\textbf{Construct-specific rules for QKD.}


\[
\text{\rm (Assign-classical)}
\quad
\Gamma
\vdash
x[i]:=a_1
\sim
x'[j]:=a_2
:
\Phi
\Longrightarrow
\ket{a_1}_{x[i]}\bra{a_1}
\land
\ket{a_2}_{x'[j]}\bra{a_2}
\land
\operatorname{proj}
\left(
\operatorname{Tr}_{x[i],x'[j]}(\Phi)
\right).
\]

\[
\text{\rm (Assign-classical-L)}
\quad
\Gamma
\vdash
x[i]:=a
\sim
\mathbf{skip}
:
\Phi
\Longrightarrow
\ket a_{x[i]}\bra a
\land
\operatorname{proj}
\left(
\operatorname{Tr}_{x[i]}(\Phi)
\right).
\]

\[
\text{\rm (Init-Quantum)}
\quad
\Gamma
\vdash
\bar q[i]:=\ket 0
\sim
\bar q'[j]:=\ket 0
:
\Phi
\Longrightarrow
\ket 0_{\bar q[i]}\bra 0
\land
\ket 0_{\bar q'[j]}\bra 0
\land
\operatorname{proj}
\left(
\operatorname{Tr}_{\bar q[i],\bar q'[j]}(\Phi)
\right).
\]

\[
\text{\rm (Init-Quantum-L)}
\quad
\Gamma
\vdash
\bar q[i]:=\ket 0
\sim
\mathbf{skip}
:
\Phi
\Longrightarrow
\ket 0_{\bar q[i]}\bra 0
\land
\operatorname{proj}
\left(
\operatorname{Tr}_{\bar q[i]}(\Phi)
\right).
\]



\[
\text{\rm (Rand)}
\quad
\frac{
\Gamma\vDash n=n'
}{
\Gamma
\vdash
z:=_{\$}\mathbf{rand}(n,\varphi)
\sim
z':=_{\$}\mathbf{rand}(n',\varphi')
:
\Phi
\land
\equiv_{(\varphi,\varphi')}
\Longrightarrow
\equiv_{(z[1,\ldots,n],z'[1,\ldots,n'])}
\land
\operatorname{proj}
\left(
\operatorname{Tr}_{z[1,\ldots,n],z'[1,\ldots,n']}(\Phi)
\right)
}.
\]

\[
\text{\rm (Rand-L)}
\quad
\Gamma
\vdash
z:=_{\$}\mathbf{rand}(n,\varphi)
\sim
\mathbf{skip}
:
\Phi
\Longrightarrow
\Pi^{\mathbb B^n}_{z[1,\ldots,n]}
\land
\operatorname{proj}
\left(
\operatorname{Tr}_{z[1,\ldots,n]}(\Phi)
\right).
\]


\[
\text{\rm (Shuffle)}
\quad
\frac{
\Gamma
\vDash
m=m'\land N=N'
}{
\Gamma
\vdash
f:=_{\$}\mathbf{shuffle}(m)_N
\sim
f':=_{\$}\mathbf{shuffle}(m')_{N'}
:
\Phi
\Longrightarrow
\operatorname{proj}
\left(
\operatorname{Tr}_{f,f'}\Phi\right)\land
\equiv_{(f,f')}
}.
\]

\[
\text{\rm (Shuffle-L)}
\quad
\Gamma
\vdash
f:=_{\$}\mathbf{shuffle}(m)_N
\sim
\mathbf{skip}
:
\Phi
\Longrightarrow
\operatorname{proj}
\left(
\operatorname{Tr}_{f}\Phi\right)\land \Pi_f^{m,N}.
\]

\[
\text{\rm (Length)}
\quad
{
\Gamma
\vdash
\iota:=\mathbf{length}(g)
\sim
\iota':=\mathbf{length}(g')
:
\Phi\land\equiv_{(g,g')}
\Longrightarrow
\operatorname{proj}\left(\operatorname{Tr}_{\iota,\iota'}\Phi\right)\land\equiv_{(\iota,\iota')}
}.
\]

\[
\text{\rm (CodeRand)}
\quad
\frac{
\Gamma
\vDash
C_m=C_m'\land n=n'
}{
\Gamma
\vdash
p :=_{\$}\textit{rand}_{n}(C_m)
\sim
p':=_{\$}\textit{rand}_{n'}(C_m')
:
\Phi
\Longrightarrow
\equiv_{(p,p')}
\land
\operatorname{proj}
\left(
\operatorname{Tr}_{p,p'}(\Phi)
\right)
}.
\]

\[
\text{\rm (CodeRand-L)}
\quad
\Gamma
\vdash
p :=_{\$}\textit{rand}_{n}(C_m)
\sim
\mathbf{skip}
:
\Phi
\Longrightarrow
\Pi_{p}
\land
\operatorname{proj}
\left(
\operatorname{Tr}_{p}(\Phi)
\right).
\]

\[
\text{\rm (IID)}
\quad
\frac{
\Gamma\vDash m=m'\land N=N'
}{
\begin{aligned}
\Gamma
\vdash\;&
f:=_{\$}\mathbf{shuffle}(m)_N;\
p:=_{\$}\mathbf{rand}(N,f)
\\
&\sim
p':=_{\$}\mathbf{rand}(N',\mathbf 1_{N'});\
f':=_{\$}\mathbf{shuffle}(m')_{N'}:
\\
&
\Phi
\Longrightarrow
\equiv_{(f,f')}
\land
\equiv_{(p[1,\ldots,N],p'[1,\ldots,N'])}
\land
\operatorname{proj}
\left(
\operatorname{Tr}_{p[1,\ldots,N],p'[1,\ldots,N']}(\Phi)
\right).
\end{aligned}
}
\]



\[
\text{\rm (UT)}
\quad
\Gamma
\vdash
\bar q[\mathbf i]\asteq U
\sim
\bar q'[\mathbf j]\asteq U'
:
\Phi
\Longrightarrow
\left(
U_{\bar q[\mathbf i]}
\otimes
U'_{\bar q'[\mathbf j]}
\right)
\Phi
\left(
U_{\bar q[\mathbf i]}
\otimes
U'_{\bar q'[\mathbf j]}
\right)^{\dagger}.
\]

\[
\text{\rm (UT-L)}
\quad
\Gamma
\vdash
\bar q[\mathbf i]\asteq U
\sim
\mathbf{skip}
:
\Phi
\Longrightarrow
U_{\bar q[\mathbf i]}
\Phi
U_{\bar q[\mathbf i]}^\dagger .
\]

We write
\[
\Gamma
\vDash
\Phi\Longrightarrow B\leftrightarrow B'
\]
to abbreviate
\[
\Gamma\vDash \Phi\land B\Longrightarrow B',
\qquad
\Gamma\vDash \Phi\land\neg B\Longrightarrow \neg B'.
\]
We write
\[
\Gamma
\vDash
\Phi\Longrightarrow_B\{\Theta_0,\Theta_1\}
\]
to abbreviate
\[
\Gamma\vDash \Phi\land B\Longrightarrow\Theta_0,
\qquad
\Gamma\vDash \Phi\land\neg B\Longrightarrow\Theta_1.
\]

\[
\text{\rm (If)}
\quad
\frac{
\Gamma
\vDash
\Phi\Longrightarrow B\leftrightarrow B'
\qquad
\Gamma
\vDash
\Phi\Longrightarrow_B\{\Theta_0,\Theta_1\}
\qquad
\Gamma
\vdash
P_0\sim P'_0:
\Theta_0\Longrightarrow\Psi
\qquad
\Gamma
\vdash
P_1\sim P'_1:
\Theta_1\Longrightarrow\Psi
}{
\Gamma
\vdash
\mathbf{if}\ B\rightarrow P_0,\ \mathbf{else}\rightarrow P_1
\sim
\mathbf{if}\ B'\rightarrow P'_0,\ \mathbf{else}\rightarrow P'_1
:
\Phi
\Longrightarrow
\Psi
}.
\]

\[
\text{\rm (If-L)}
\quad
\frac{
\Gamma
\vDash
\Phi\Longrightarrow_B\{\Theta_0,\Theta_1\}
\qquad
\Gamma
\vdash
P_0\sim P:
\Theta_0\Longrightarrow\Psi
\qquad
\Gamma
\vdash
P_1\sim P:
\Theta_1\Longrightarrow\Psi
}{
\Gamma
\vdash
\mathbf{if}\ B\rightarrow P_0,\ \mathbf{else}\rightarrow P_1
\sim
P
:
\Phi
\Longrightarrow
\Psi
}.
\]




\begin{definition}[Measurement Condition]
Let
\[
\mathcal M_1=\{M_1^{r}\}_{r\in[m]},
\qquad
\mathcal M_2=\{M_2^{r}\}_{r\in[m]}
\]
be two measurements with the same outcome set
\([m]=\{1,\ldots,m\}\). The measurement condition
\[
\Gamma
\vDash
\mathcal M_1\approx\mathcal M_2:
\Phi_{\mathcal V}
\Longrightarrow
\Psi_{\mathcal W}
\]
means that, for every \(r\in[m]\) and every pair of well-typed
c-q states \(\Delta,\Delta'\), if
\[
\Delta\ \Phi_{\mathcal V}\ \Delta'
\]
has a witness \(\pi\) satisfying
\[
\operatorname{supp}_{\mathsf{cl}}(\pi)\models\Gamma,
\]
then
\[
M_1^r\Delta M_1^{r\dagger}
\ \Psi_{\mathcal W}\
M_2^r\Delta'M_2^{r\dagger}.
\]

For a one-sided measurement, the condition
\[
\Gamma
\vDash
\mathcal M\approx\mathcal I:
\Phi_{\mathcal V}
\Longrightarrow
\Psi_{\mathcal W}
\]
means that, for every pair of well-typed c-q states
\(\Delta,\Delta'\), if
\[
\Delta\ \Phi_{\mathcal V}\ \Delta'
\]
has a witness \(\pi\) satisfying
\[
\operatorname{supp}_{\mathsf{cl}}(\pi)\models\Gamma,
\]
then there exists a family of c-q states
\[
\{\Delta'_r\}_{r\in[m]}
\qquad\text{such that}\qquad
\sum_{r\in[m]}\Delta'_r=\Delta',
\]
and, for every \(r\in[m]\),
\[
M^r\Delta M^{r\dagger}
\ \Psi_{\mathcal W}\
\Delta'_r.
\]
\end{definition}

For a classical logical reference \(z\) with values in \([m]\), define
\[
\Pi_z^{[m]}
\stackrel{\triangle}{=}
\sum_{r\in[m]}
\ket r_z\bra r .
\]
For paired classical logical references \(z,z'\), define
\[
\equiv^{[m]}_{(z,z')}
\stackrel{\triangle}{=}
\sum_{r\in[m]}
\ket r_z\bra r
\otimes
\ket r_{z'}\bra r .
\]

\[
\text{\rm (Meas)}
\quad
\frac{
\Gamma
\vDash
\mathcal M_1\approx\mathcal M_2:
\Phi_{\mathcal V}
\Longrightarrow
\Psi_{\mathcal W}
}{
\Gamma
\vdash
x[i]:=\mathcal M_1[\bar q[j]]
\sim
x'[i']:=\mathcal M_2[\bar q'[j']]
:
\Phi_{\mathcal V}
\Longrightarrow
\equiv^{[m]}_{(x[i],x'[i'])}
\land
\operatorname{proj}
\left(
\operatorname{Tr}_{x[i],x'[i']}(\Psi_{\mathcal W})
\right)
}.
\]

\[
\text{\rm (Meas-L)}
\quad
\frac{
\Gamma
\vDash
\mathcal M\approx\mathcal I:
\Phi_{\mathcal V}
\Longrightarrow
\Psi_{\mathcal W}
}{
\Gamma
\vdash
x[i]:=\mathcal M[\bar q[j]]
\sim
\mathbf{skip}
:
\Phi_{\mathcal V}
\Longrightarrow
\Pi^{[m]}_{x[i]}
\land
\operatorname{proj}
\left(
\operatorname{Tr}_{x[i]}(\Psi_{\mathcal W})
\right)
}.
\]

\[
\text{\rm (CNOT)}
\quad
{\begin{aligned}
\Gamma
&\vdash
\bar q_1,\bar q_2\asteq \mathrm{CNOT};
\ c[t]:=\mathcal M_{\mathrm{com}}[\bar q_1]
\sim
c'[t']:=\mathcal M_{\mathrm{com}}[\bar q'_1];
\ \mathbf{if}\ c'[t']=0\rightarrow\mathbf{skip},
\ \mathbf{else}\rightarrow \bar q'_2\asteq X :\\ 
&
\equiv_{((\bar q_1,\bar q_2),(\bar q'_1,\bar q'_2))}
\Longrightarrow
\equiv_{(c[t],c'[t'])}
\land
\equiv_{((\bar q_1,\bar q_2),(\bar q'_1,\bar q'_2))}
\end{aligned}}
\]

\[
\text{\rm (If-Stru-1)}
\quad
\frac{
\Gamma\vdash
\operatorname{upd}(c_2)
\#
\operatorname{acc}(b)
}{
\Gamma
\vdash
c_2;\
\mathbf{if}\ b=0\rightarrow c_0,\ \mathbf{else}\rightarrow c_1
\sim
\mathbf{if}\ b=0\rightarrow c_2;c_0,\ \mathbf{else}\rightarrow c_2;c_1
:
\equiv_{\mathcal V}
\Longrightarrow
\equiv_{\mathcal V}
}.
\]

\[
\text{\rm (If-Stru-2)}
\quad
\Gamma
\vdash
(\mathbf{if}\ b=0\rightarrow c_0,\ \mathbf{else}\rightarrow c_1);c_2
\sim
\mathbf{if}\ b=0\rightarrow c_0;c_2,\ \mathbf{else}\rightarrow c_1;c_2
:
\equiv_{\mathcal V}
\Longrightarrow
\equiv_{\mathcal V}.
\]

\begin{definition}[Bad Event Assertion]
Let \(bad\) be a bit-valued failure-event view, also written
\[
bad
\stackrel{\triangle}{=}
\ket 1_{bad}\bra 1
\]
for the projector where \(bad=1\), and define
\[
\neg bad
\stackrel{\triangle}{=}
I-bad
=
\ket 0_{bad}\bra 0.
\]
For a left-side bad event write
\[
\neg bad\otimes I
\]
to act on the left and as identity on the right. This proof-side flag is not a
terminal QKD observation component.
\end{definition}

\[
\text{\rm (UTB)}
\quad
\frac{
\Gamma
\vdash
P'_1\sim_{\delta'}P:
\equiv_{\mathcal V}
\Longrightarrow
\equiv_{\mathcal V}\land(\neg bad\otimes I)
\qquad
\Gamma
\vdash
P\sim_{\delta}P':
\equiv_{\mathcal V}
\Longrightarrow
\equiv_{\mathcal V}
}{
\Gamma
\vdash
P'_1\sim_{\delta+\delta'}P':
\equiv_{\mathcal V}
\Longrightarrow
\equiv_{\mathcal V}\land(\neg bad\otimes I)
}.
\]

\begin{definition}[Uniformity in a Basis]
Let
\(\mathcal B=\{\ket{i}\}_{i=1}^{D}\)
be an orthonormal basis. \(\rho\) is uniform in \(\mathcal B\) iff measurement
in \(\mathcal B\) is uniform:
\[
p_i=\operatorname{Tr}\left(\ket{i}\bra{i}\rho\right)
=\langle i|\rho|i\rangle
=\frac{1}{D}\qquad
(1\le i\le D).
\]
\end{definition}

Uniformity is observable, not uniquely projective; here it is used through the
sufficient condition \(\ket{+}^{\otimes\ell}\) before computational-basis
measurement. For
\[
\mathcal B=
\{\ket{s}_{\bar q}:s\in\{0,1\}^{\ell}\},
\]
define
\[
\ket +_{\bar q}
=
\frac{1}{\sqrt{2^\ell}}
\sum_{s\in\{0,1\}^{\ell}}
\ket s_{\bar q},
\qquad
\Phi_{\bar q}^{\mathcal B}
=
\ket +_{\bar q}\bra +.
\]

\[
\text{\rm (Uniform)}
\quad
\frac{
\Gamma
\vDash
\Phi
\Longrightarrow
\Phi_{\bar q}^{\mathcal B}
}{
\Gamma
\vdash
p
:=
\mathcal M_{\mathcal B}[\bar q[1,\ldots,\ell]]
\sim
p'
:=_{\$}
\mathbf{rand}(\ell)
:
\Phi
\Longrightarrow
\equiv_{(p,p')}
}.
\]

% ============================================================
% Descriptor-sensitive sift
% ============================================================
\begin{definition}[Descriptor Alias Notation]
Let \(D_f\) be a descriptor command controlled by the filter \(f\). For a view
set \(\mathcal V\), write
\[
\operatorname{src}_{D_f}(\mathcal V),
\qquad
\operatorname{tgt}_{D_f}(\mathcal V)
\]
for the source-side and target-side versions of \(\mathcal V\).

For
\[
D_f
=
(g_0,g_1)
:=
\mathbf{select}_{f}^{N}\ \mathbf{from}~J~\mathbf{for}\ g,
\]
\[
\operatorname{src}_{D_f}(g)=g,
\qquad
\operatorname{tgt}_{D_f}(g)=(g_0,g_1).
\]

If
\[
D_f
=
(\mathcal P_0,\mathcal P_1)
:=
\mathbf{select}_{f}^{N}\ \mathbf{from}~J~\mathbf{for}\ \mathcal P,
\]
then
\[
\operatorname{src}_{D_f}(\mathcal P)=\mathcal P,
\qquad
\operatorname{tgt}_{D_f}(\mathcal P)=(\mathcal P_0,\mathcal P_1).
\]

If \(D_f\) selects \(b,\bar q\) into \(b_c,\bar q_c,b_k,\bar q_k\), and a
register command produces \(x_c,x_k\), the rule instance may use
\[
\operatorname{tgt}_{D_f}(b,x)=((b_c,x_c),(b_k,x_k)).
\]

For
\[
D_f
=
g
:=
\mathbf{merge}_{f}^{N}\ \mathbf{from}~J~\mathbf{for}\ g_0,g_1,
\]
\[
\operatorname{src}_{D_f}(g)=(g_0,g_1),
\qquad
\operatorname{tgt}_{D_f}(g)=g.
\]

If
\[
D_f
=
\mathcal P
:=
\mathbf{merge}_{f}^{N}\ \mathbf{from}~J~\mathbf{for}\ \mathcal P_0,\mathcal P_1,
\]
then
\[
\operatorname{src}_{D_f}(\mathcal P)=(\mathcal P_0,\mathcal P_1),
\qquad
\operatorname{tgt}_{D_f}(\mathcal P)=\mathcal P.
\]
\end{definition}

\begin{definition}[Alias Compatibility Check]
The proof-side check
\[
\Gamma
\vdash
D_f[\mathcal P]
\cong
D'_{f'}[\mathcal P']
\]
means that the two displayed descriptor commands have the same filter-layout
form.

The filter relation is written
\[
\equiv_{(f,f')}.
\]

The role of
\[
\Gamma
\vdash
D_f[\mathcal P]
\cong
D'_{f'}[\mathcal P']
\]
is that the two constructions have the same descriptor shape: both are
\textbf{select} constructions or both are \textbf{merge} constructions, with
the same displayed length parameter \(N\), the same source-length expression
\(J\), and compatible corresponding components.
\end{definition}

\begin{definition}[Fiber Side Condition]
The side condition
\[
\Gamma
\vDash
C_L
\sim_{\mathsf{fiber}}
C_R
:
\equiv_{\mathcal U}
\Longrightarrow
\equiv_{\mathcal W}
\]
is a certified local adapter for register commands that are moved across a
nearby descriptor construction such as select or merge.

Its purpose is to justify that, after a descriptor command changes the logical
view names, the same register-side operation can be written using the adapted
source or target view names.

The commands \(C_L,C_R\) must be descriptor-preserving:
\[
\operatorname{dupd}(C_L)=\operatorname{dupd}(C_R)=\emptyset.
\]

The input view set \(\mathcal U\) and output view set \(\mathcal W\) are
determined by the displayed register schemas. They are not inferred from
arbitrary descriptor commands.

Their register effects are recorded by
\[
\operatorname{racc}(C_L),
\operatorname{rupd}(C_L),
\qquad
\operatorname{racc}(C_R),
\operatorname{rupd}(C_R).
\]
\end{definition}



\[
\text{\rm (Select)}
\quad
\frac{
\Gamma\vDash J=J'
}{
\begin{aligned}
\Gamma
\vdash\;&
(\mathcal P_0,\mathcal P_1)
:=
\mathbf{select}_{f}^{N}\ \mathbf{from}\ J\ \mathbf{for}\ \mathcal P
\sim
(\mathcal P'_0,\mathcal P'_1)
:=
\mathbf{select}_{f'}^{N}\ \mathbf{from}\ J'\ \mathbf{for}\ \mathcal P'
:
\\
&
\equiv_{(f,f')}\land\equiv_{(\mathcal P,\mathcal P')}
\Longrightarrow
\equiv_{((\mathcal P_0,\mathcal P_1),(\mathcal P'_0,\mathcal P'_1))}
\end{aligned}
}
\]

\[
\text{\rm (Merge)}
\quad
\frac{
\Gamma\vDash J=J'
}{
\begin{aligned}
\Gamma
\vdash\;&
\mathcal P
:=
\mathbf{merge}_{f}^{N}\ \mathbf{from}\ J\ \mathbf{for}\ \mathcal P_0,\mathcal P_1
\sim
\mathcal P'
:=
\mathbf{merge}_{f'}^{N}\ \mathbf{from}\ J'\ \mathbf{for}\ \mathcal P'_0,\mathcal P'_1
:
\\
&
\equiv_{(f,f')}\land\equiv_{((\mathcal P_0,\mathcal P_1),
          (\mathcal P'_0,\mathcal P'_1))}
\Longrightarrow
\equiv_{(\mathcal P,\mathcal P')}
\end{aligned}
}
\]

\[
\text{\rm (Select-L)}
\quad{
\begin{aligned}
\Gamma
\vdash\;&
(\mathcal P_0,\mathcal P_1)
:=
\mathbf{select}_{f}^{N}\ \mathbf{from}\ J\ \mathbf{for}\ \mathcal P
\sim
\mathbf{skip}
:
\equiv_{(\mathcal P^f,\mathcal U')}
\Longrightarrow
\equiv_{((\mathcal P_0,\mathcal P_1),\mathcal U')}.
\end{aligned}
}
\]

\[
\text{\rm (Merge-L)}
\quad
{
\begin{aligned}
\Gamma
\vdash\;&
\mathcal P
:=
\mathbf{merge}_{f}^{N}\ \mathbf{from}\ J\ \mathbf{for}\ \mathcal P_0,\mathcal P_1
\sim
\mathbf{skip}
:
\equiv_{((\mathcal P_0,\mathcal P_1)^{\uparrow f},\mathcal U')}
\Longrightarrow
\equiv_{(\mathcal P,\mathcal U')}.
\end{aligned}
}
\]

\[
\text{\rm (Sift)}
\quad
\frac{
\begin{array}{c}
\Gamma
\vdash
D_f[\mathcal P_{\mathsf{pre}}]
\cong
D'_{f'}[\mathcal P'_{\mathsf{post}}]
\\[0.6ex]
\Gamma
\vDash
C_{\operatorname{tgt}_{D_f}}
\sim_{\mathsf{fiber}}
C_{\operatorname{src}_{D'_{f'}}}
:
\equiv_{
\left(
\operatorname{tgt}_{D_f}(\mathcal V_{\mathsf{pre}}),
\operatorname{src}_{D'_{f'}}(\mathcal V'_{\mathsf{pre}})
\right)
}
\Longrightarrow
\equiv_{
\left(
\operatorname{tgt}_{D_f}(\mathcal V_{\mathsf{post}}),
\operatorname{src}_{D'_{f'}}(\mathcal V'_{\mathsf{post}})
\right)
}
\end{array}
}{
\begin{aligned}
\Gamma
\vdash\;&
D_f[\mathcal P_{\mathsf{pre}}];
C_{\operatorname{tgt}_{D_f}}
\sim
C_{\operatorname{src}_{D'_{f'}}};
D'_{f'}[\mathcal P'_{\mathsf{post}}]:
\equiv_{(f,f')}
\land
\equiv_{
\left(
\operatorname{src}_{D_f}(\mathcal V_{\mathsf{pre}}),
\operatorname{src}_{D'_{f'}}(\mathcal V'_{\mathsf{pre}})
\right)
}
\Longrightarrow
\equiv_{
\left(
\operatorname{tgt}_{D_f}(\mathcal V_{\mathsf{post}}),
\operatorname{tgt}_{D'_{f'}}(\mathcal V'_{\mathsf{post}})
\right)
}.
\end{aligned}
}
\]



\paragraph{Certified fiber instances.}

For basis readout with matching descriptors:
\[
\Gamma
\vDash
\mathsf{Read}_{b}^{E}(x;\bar q)
\sim_{\mathsf{fiber}}
\mathsf{Read}_{b'}^{E}(x';\bar q')
:
\equiv_{((b,\bar q),(b',\bar q'))}
\Longrightarrow
\equiv_{((b,x),(b',x'))}.
\]

For split-vs-unsplit basis readout, where
\(m_0=\operatorname{wt}_{E}(\neg f)\) and
\(m_1=\operatorname{wt}_{E}(f)\):
\[
\Gamma
\vDash
\mathsf{Read}_{b_0}^{m_0}(x_0;\bar q_0);
\mathsf{Read}_{b_1}^{m_1}(x_1;\bar q_1)
\sim_{\mathsf{fiber}}
\mathsf{Read}_{b'}^{E}(x';\bar q')
:
\equiv_{(((b_0,\bar q_0),(b_1,\bar q_1)),(b',\bar q'))}
\Longrightarrow
\equiv_{(((b_0,x_0),(b_1,x_1)),(b',x'))}.
\]

For computational-basis bit preparation with matching descriptors:
\[
\Gamma
\vDash
\mathsf{PrepBits}^{E}(a;\bar q)
\sim_{\mathsf{fiber}}
\mathsf{PrepBits}^{E}(a';\bar q')
:
\equiv_{(a,a')}
\Longrightarrow
\equiv_{((a,\bar q),(a',\bar q'))}.
\]

For split-vs-unsplit computational-basis bit preparation:
\[
\Gamma
\vDash
\mathsf{PrepBits}^{E}(a;\bar q)
\sim_{\mathsf{fiber}}
\mathsf{PrepBits}^{m_0}(a'_0;\bar q'_0);
\mathsf{PrepBits}^{m_1}(a'_1;\bar q'_1)
:
\equiv_{(a,(a'_0,a'_1))}
\Longrightarrow
\equiv_{((a,\bar q),((a'_0,\bar q'_0),(a'_1,\bar q'_1)))}.
\]

For basis-indexed bit preparation with matching descriptors:
\[
\Gamma
\vDash
\mathsf{PrepBits}^{E}_{b}(a;\bar q)
\sim_{\mathsf{fiber}}
\mathsf{PrepBits}^{E}_{b'}(a';\bar q')
:
\equiv_{((a,b),(a',b'))}
\Longrightarrow
\equiv_{((a,b,\bar q),(a',b',\bar q'))}.
\]

For split-vs-unsplit basis-indexed bit preparation:
\[
\Gamma
\vDash
\mathsf{PrepBits}^{E}_{b}(a;\bar q)
\sim_{\mathsf{fiber}}
\mathsf{PrepBits}^{m_0}_{b'_0}(a'_0;\bar q'_0);
\mathsf{PrepBits}^{m_1}_{b'_1}(a'_1;\bar q'_1)
:
\equiv_{((a,b),((a'_0,b'_0),(a'_1,b'_1)))}
\Longrightarrow
\equiv_{
((a,b,\bar q),((a'_0,b'_0,\bar q'_0),(a'_1,b'_1,\bar q'_1)))
}.
\]

\paragraph{Derived sift-measure instance.}

Let
\[
m_c=\operatorname{wt}_{J}(\neg f),
\qquad
m_k=\operatorname{wt}_{J}(f).
\]
Taking
\[
D_f[\mathcal P_{\mathsf{pre}}]
=
\mathbf{select}_{f}^{N}\ \mathbf{from}\ J\ \mathbf{for}\ b,\bar q,
\qquad
D'_{f'}[\mathcal P'_{\mathsf{post}}]
=
\mathbf{select}_{f'}^{N}\ \mathbf{from}\ J\ \mathbf{for}\ b',x',
\]
\[
\mathcal V_{\mathsf{pre}}=(b,\bar q),
\qquad
\mathcal V_{\mathsf{post}}=(b,x),
\]
gives:
\[
\begin{aligned}
\Gamma
\vdash\;&
(b_c,\bar q_c),(b_k,\bar q_k)
:=
\mathbf{select}_{f}^{N}\ \mathbf{from}\ J\ \mathbf{for}\ b,\bar q;
\\
&
\mathsf{Read}_{b_c}^{m_c}(x_c;\bar q_c);
\mathsf{Read}_{b_k}^{m_k}(x_k;\bar q_k)
\\
&\sim
\mathsf{Read}_{b'}^{J}(x';\bar q');
\\
&
(b'_c,x'_c),(b'_k,x'_k)
:=
\mathbf{select}_{f'}^{N}\ \mathbf{from}\ J\ \mathbf{for}\ b',x'
\\
&:
\equiv_{(f,f')}
\land
\equiv_{(b,b')}
\land
\equiv_{(\bar q,\bar q')}
\\
&\Longrightarrow
\equiv_{\bigl(
(b_c,x_c,b_k,x_k),
(b'_c,x'_c,b'_k,x'_k)
\bigr)} .
\end{aligned}
\]


\paragraph{Derived sift-prepare instance.}

Let
\[
m_c=\operatorname{wt}_{J}(\neg f),
\qquad
m_k=\operatorname{wt}_{J}(f).
\]

For computational-basis preparation, taking
\[
D_f[\mathcal P_{\mathsf{pre}}]
=
\mathbf{merge}_{f}^{N}\ \mathbf{from}\ J\
\mathbf{for}\ (a_c,\bar q_c),(a_k,\bar q_k),
\]
\[
D'_{f'}[\mathcal P'_{\mathsf{post}}]
=
\mathbf{merge}_{f'}^{N}\ \mathbf{from}\ J\
\mathbf{for}\ (a'_c,\bar q'_c),(a'_k,\bar q'_k),
\]
\[
\mathcal V_{\mathsf{pre}}=a,
\qquad
\mathcal V_{\mathsf{post}}=(a,\bar q),
\]
gives:
\[
\begin{aligned}
\Gamma
\vdash\;&
(a,\bar q)
:=
\mathbf{merge}_{f}^{N}\ \mathbf{from}\ J\
\mathbf{for}\ (a_c,\bar q_c),(a_k,\bar q_k);
\\
&
\mathsf{PrepBits}^{J}(a;\bar q)
\\
&\sim
\mathsf{PrepBits}^{m_c}(a'_c;\bar q'_c);
\mathsf{PrepBits}^{m_k}(a'_k;\bar q'_k);
\\
&
(a',\bar q')
:=
\mathbf{merge}_{f'}^{N}\ \mathbf{from}\ J\
\mathbf{for}\ (a'_c,\bar q'_c),(a'_k,\bar q'_k)
\\
&:
\equiv_{(f,f')}
\land
\equiv_{(a_c,a'_c)}
\land
\equiv_{(a_k,a'_k)}
\\
&\Longrightarrow
\equiv_{(a,a')}
\land
\equiv_{(\bar q,\bar q')} .
\end{aligned}
\]

For basis-indexed preparation, taking
\[
D_f[\mathcal P_{\mathsf{pre}}]
=
\mathbf{merge}_{f}^{N}\ \mathbf{from}\ J\
\mathbf{for}\
(a_c,b_c,\bar q_c),(a_k,b_k,\bar q_k),
\]
\[
D'_{f'}[\mathcal P'_{\mathsf{post}}]
=
\mathbf{merge}_{f'}^{N}\ \mathbf{from}\ J\
\mathbf{for}\
(a'_c,b'_c,\bar q'_c),(a'_k,b'_k,\bar q'_k),
\]
\[
\mathcal V_{\mathsf{pre}}=(a,b),
\qquad
\mathcal V_{\mathsf{post}}=(a,b,\bar q),
\]
gives:
\[
\begin{aligned}
\Gamma
\vdash\;&
(a,b,\bar q)
:=
\mathbf{merge}_{f}^{N}\ \mathbf{from}\ J\
\mathbf{for}\ (a_c,b_c,\bar q_c),(a_k,b_k,\bar q_k);
\\
&
\mathsf{PrepBits}^{J}_{b}(a;\bar q)
\\
&\sim
\mathsf{PrepBits}^{m_c}_{b'_c}(a'_c;\bar q'_c);
\mathsf{PrepBits}^{m_k}_{b'_k}(a'_k;\bar q'_k);
\\
&
(a',b',\bar q')
:=
\mathbf{merge}_{f'}^{N}\ \mathbf{from}\ J\
\mathbf{for}\ (a'_c,b'_c,\bar q'_c),(a'_k,b'_k,\bar q'_k)
\\
&:
\equiv_{(f,f')}
\land
\equiv_{((a_c,b_c),(a'_c,b'_c))}
\land
\equiv_{((a_k,b_k),(a'_k,b'_k))}
\\
&\Longrightarrow
\equiv_{(a,a')}
\land
\equiv_{(b,b')}
\land
\equiv_{(\bar q,\bar q')} .
\end{aligned}
\]


\section{Syntax Sugar}
\label{app:syntax-sugar}

\begin{itemize}

    \item \textbf{Initialization of quantum variables:}
    \[
    \bar q[i]:=|1\rangle\stackrel{\triangle}{=}
    \bar q[i]:=|0\rangle;\ \bar q[i]\asteq X;
    \]
    
    \[
    \overline{q} := \ket{0}^{\otimes n}
    \stackrel{\triangle}{=}
    \bar{q}[1] := \ket{0};\ \cdots;\ \bar{q}[n] := \ket{0};\]

    \[
    \overline{q} := \ket{1}^{\otimes n}
    \stackrel{\triangle}{=}
    \bar{q}[1] := \ket{1};\ \cdots;\ \bar{q}[n] := \ket{1}.\]

    \item {\textbf{Dummy channel-order syntax sugar.}}
    
    {Private auxiliary channel-side systems must already belong
    to the fixed channel/adversarial workspace used with}
    \[
    {
    U_{\textbf{channel}}.
    }
    \]
    {They are ordinary systems outside the retained terminal QKD
    observation.}
    
    \item \textbf{Measurement in the computational basis:}
    For a qubit register
    \(\bar q=(\bar q[1],\ldots,\bar q[n])\)
    and a bit-valued classical register
    \(x=(x[1],\ldots,x[n])\), we write
    \[
    x[1,\ldots,n]
    :=
    \mathcal M_{\mathrm{com}}[\bar q[1,\ldots,n]]\]
    as syntactic sugar for
    \(x[1]:=\mathcal M_{\mathrm{com}}[\bar q[1]];
    \ \cdots;\
    x[n]:=\mathcal M_{\mathrm{com}}[\bar q[n]],\)
    where the single-qubit computational-basis measurement is
    \(\mathcal M_{\mathrm{com}}
    \stackrel{\triangle}{=}
    \{\ket 0\bra 0,\ket 1\bra 1\}.\)

    \item \(\mathsf{Read}_{b}^{E}(x;\bar q)\):
    For a bit-valued basis view \(b\), a bit-valued output view \(x\), and a
    quantum view \(\bar q\), define
    \[
    \mathsf{Read}_{b}^{E}(x;\bar q)
    \]
    as syntax sugar for
    \[
    \mathbf{for}^{E}\ i=1\ \mathbf{to}\ E
    \quad
    \mathbf{if}\ b[i]=0
    \rightarrow
    x[i]:=\mathcal M_{\mathrm{com}}[\bar q[i]],
    \]
    \[
    \hspace{27mm}
    \mathbf{else}
    \rightarrow
    x[i]:=\mathcal M_{\{+,-\}}[\bar q[i]]
    \quad
    \mathbf{end}.
    \]
    It is a descriptor-preserving register operation.

    \item \(\mathsf{PrepBits}^{E}(a;\bar q)\) and
    \(\mathsf{PrepBits}^{E}_{b}(a;\bar q)\):

    For a bit-valued input view \(a\), a bit-valued basis view \(b\), and a
    quantum view \(\bar q\), define basis-indexed preparation
    \[
    \mathsf{PrepBits}^{E}_{b}(a;\bar q)
    \]
    as syntax sugar for
    \[
    \mathbf{for}^{E}\ i=1\ \mathbf{to}\ E
    \quad
    \mathbf{if}\ b[i]=0
    \rightarrow S_0(i),
    \quad
    \mathbf{else}
    \rightarrow S_1(i)
    \quad
    \mathbf{end},
    \]
    where
    \[
    S_0(i)
    \stackrel{\triangle}{=}
    \mathbf{if}\ a[i]=0
    \rightarrow
    \bar q[i]:=\ket0,
    \quad
    \mathbf{else}
    \rightarrow
    \bar q[i]:=\ket1,
    \]
    and
    \[
    S_1(i)
    \stackrel{\triangle}{=}
    \mathbf{if}\ a[i]=0
    \rightarrow
    \bar q[i]:=\ket+,
    \quad
    \mathbf{else}
    \rightarrow
    \bar q[i]:=\ket- .
    \]

    Omitting the basis subscript means computational-basis preparation:
    \[
    \mathsf{PrepBits}^{E}(a;\bar q)
    \stackrel{\triangle}{=}
    \mathsf{PrepBits}^{E}_{\mathbf 0_E}(a;\bar q),
    \]
    where
    \[
    \mathbf 0_E[i]=0
    \qquad
    (1\le i\le E).
    \]

    Thus
    \[
    \mathsf{PrepBits}^{E}(a;\bar q)
    \]
    expands to
    \[
    \mathbf{for}^{E}\ i=1\ \mathbf{to}\ E
    \quad
    \mathbf{if}\ a[i]=0
    \rightarrow
    \bar q[i]:=\ket0,
    \quad
    \mathbf{else}
    \rightarrow
    \bar q[i]:=\ket1
    \quad
    \mathbf{end}.
    \]

    Both forms are descriptor-preserving register operations. Their output
    quantum state is determined by the displayed classical controls and not by
    the previous state of \(\bar q[i]\).

    \item \textbf{Classically indexed unitary operations on a fixed register:}
    For a unitary operator \(U\) acting on \(k\) qubits, we write
    \(\bar q[i_1,\ldots,i_k]\asteq U\)
    for applying \(U\) to ordered variables
    \(\bar q[i_1],\ldots,\bar q[i_k]\); \(k=1\) is primitive.
    Let \(U\stackrel{\triangle}{=}(U_i)_{1\le i\le K}\) be a finite family
    of unitary operators on the Hilbert space associated with \(\bar q\),
    with \(K\ge 1\), and let \(e\) be an \(\mathsf{Int}\)-typed expression.
    We write
    \[
    \bar q\asteq U(e)
    \]
    for the finite case expansion
    \[
    \textbf{if } e<1\lor e>K \rightarrow \textbf{abort},
    \ \textbf{else}\rightarrow S_1;\cdots;S_K,
    \]
    where, for each \(1\le i\le K\),
    \[
    S_i
    \stackrel{\triangle}{=}
    \textbf{if } e=i \rightarrow \bar q\asteq U_i,
    \ \textbf{else}\rightarrow\mathbf{skip}.
    \]
    Since \(S_i\) does not change \(e\), exactly one valid branch applies
    \(U_i\); out-of-range aborts.

    \item \textbf{Unitary operations on selected qubits:}
    Let \(1\le k\le n\), and let \(e,e_1,\ldots,e_k\) be
    \(\mathsf{Int}\)-typed expressions. The statement
    \[
    \bar q[e_1,\ldots,e_k]\asteq U(e),
    \]
    where \(U=(U_i)_{1\le i\le K}\) is as in the previous clause and each
    \(U_i\) acts on \(k\) qubits, applies \(U_i\) on the quantum systems
    \(\bar q[i_1],\ldots,\bar q[i_k]\) whenever \(e\) evaluates to \(i\)
    and each \(e_j\) evaluates to a distinct valid \(i_j\). The guard
    \[
    \exists\,1\le r\le k.\,(e_r<1\lor e_r>n)
    \ \lor\
    \exists\,1\le r<s\le k.\,(e_r=e_s)
    \]
    abbreviates the corresponding finite predicate. The statement denotes
    \[
    \textbf{if } B_{\mathrm{bad}}\rightarrow \textbf{abort},
    \ \textbf{else}\rightarrow S_1;S_2;\cdots;S_{|R|},
    \]
    where
    \(B_{\mathrm{bad}}
    \stackrel{\triangle}{=}
    \bigvee_{r=1}^{k}(e_r<1\lor e_r>n)
    \ \lor\
    \bigvee_{1\le r<s\le k}(e_r=e_s),\)
    and
    \(R
    \stackrel{\triangle}{=}
    \{(i_1,\ldots,i_k):
    1\le i_j\le n\ \text{ for all }j,\ \text{and }i_1,\ldots,i_k
    \text{ are distinct}\}.\)
    Fix an enumeration
    \[
    R=\{(i_1^\ell,\ldots,i_k^\ell):1\le \ell\le |R|\}.
    \]
    For each \(1\le \ell\le |R|\), define
    \[
    S_\ell
    \stackrel{\triangle}{=}
    \textbf{if }e_1=i_1^\ell\land\cdots\land e_k=i_k^\ell
    \rightarrow
    \bar q[i_1^\ell,\ldots,i_k^\ell]\asteq U(e),
    \ \textbf{else}\rightarrow\mathbf{skip}.
    \]

    \item \textbf{Preparation and measurement in the \(\{+,-\}\) basis:}
    For one qubit:
    \begin{align*}
        \bar q[i] := \ket{+}
        &\stackrel{\triangle}{=}
        \bar q[i] := \ket{0};\ \bar q[i]\asteq H,\\
        \bar q[i] := \ket{-}
        &\stackrel{\triangle}{=}
        \bar q[i] := \ket{0};\ \bar q[i]\asteq X;\ \bar q[i]\asteq H,\\
        c[t] := \mathcal M_{\{+,-\}}[\bar q[i]]
        &\stackrel{\triangle}{=}
        \bar q[i]\asteq H;\ c[t]:=\mathcal M_{\mathrm{com}}[\bar q[i]].
    \end{align*}

    \item \textbf{Boolean XOR notation:}
    For bit-valued expressions \(E_1,E_2\), we may use
    \[
    E_1\oplus E_2=0
    \quad\stackrel{\triangle}{=}\quad
    E_1=E_2,
    \qquad\text{and}\qquad
    E_1\oplus E_2=1
    \quad\stackrel{\triangle}{=}\quad
    E_1\neq E_2.
    \]
    Assignment notation
    \[
    x[t] := E_1\oplus E_2
    \]
    as syntactic sugar for
    \[
    \textbf{if } E_1=E_2 \rightarrow x[t]:=0,
    \ \textbf{else}\rightarrow x[t]:=1.
    \]

    \item \textbf{Fixed finite classical-map syntax sugars.}

    Let
    \[
    F:D_1\times\cdots\times D_s\to D
    \]
    be a fixed external classical map over finite protocol domains. The command
    \[
    y:=F(x_1,\ldots,x_s)
    \]
    is finite case analysis over \(D_1\times\cdots\times D_s\), using only
    conditionals and constant assignments. For each tuple
    \[
    d=(d_1,\ldots,d_s)\in D_1\times\cdots\times D_s,
    \]
    let
    \[
    c_d=F(d_1,\ldots,d_s).
    \]
    Then \(y:=F(x_1,\ldots,x_s)\) expands into the finite sequence of guarded
    assignments
    \[
    \prod_{d\in D_1\times\cdots\times D_s}
    \left(
    \mathbf{if}\ x_1=d_1\land\cdots\land x_s=d_s
    \rightarrow y:=c_d,
    \ \mathbf{else}\rightarrow\mathbf{skip}
    \right),
    \]
    for any fixed enumeration of the finite domain.

    Vector-valued fixed maps are expanded componentwise. Thus commands such as
    \[
    k_b:=\mathsf{qecBob}(k',x-\nu_k),
    \qquad
    k_a:=\mathsf{qecAlice}(k,x-\nu_k),
    \qquad
    \hat e_X:=D_X(s_Z),
    \qquad
    \hat e_Z:=D_Z(s_X)
    \]
    are finite classical-map syntax sugars, not primitives.

    \item \textbf{CNOT gate:}
    For \(k\ge 1\), we write
    \(\bar q[e_1,\ldots,e_k,e_{k+1},\ldots,e_{2k}]\asteq \mathrm{CNOT}\)
    for the parallel application of \(k\) CNOT gates, where \(\bar q[e_j]\)
    is the control and \(\bar q[e_{k+j}]\) is the target of the \(j\)-th CNOT
    gate:
    \[
    \bar q[e_1,\ldots,e_k,e_{k+1},\ldots,e_{2k}]\asteq \mathrm{CNOT}
    \stackrel{\triangle}{=}
    \bar q[e_1,e_{k+1},e_2,e_{k+2},\ldots,e_k,e_{2k}]
    \asteq \mathrm{CNOT}^{\otimes k}.
    \]
    Range and distinctness come from selected-qubit unitary sugar. If
    \(\bar q=(\bar q[1],\bar q[2])\) is a two-qubit register, then
    \(\bar q\asteq \mathrm{CNOT}
    \quad\text{stands for}\quad
    \bar q[1,2]\asteq \mathrm{CNOT},\)
    with \(\bar q[1]\) as control and \(\bar q[2]\) as target.

    \item \textbf{Pauli-parity and CSS-syndrome measurement.}
    
    Matrices here are fixed external data, not sampled program variables.
    
    Let \(h=(h_1,\ldots,h_N)\in\{0,1\}^{N}\) be a fixed binary protocol vector.
    Define
    \[
    \operatorname{supp}(h)
    \stackrel{\triangle}{=}
    \{t: h_t=1\}.
    \]
    For a quantum block
    \[
    \bar r=(\bar r[1],\ldots,\bar r[N]),
    \]
    define
    \[
    Z(h)\stackrel{\triangle}{=}\bigotimes_{t=1}^{N}Z_t^{h_t},
    \qquad
    X(h)\stackrel{\triangle}{=}\bigotimes_{t=1}^{N}X_t^{h_t}.
    \]
    The syndrome bit convention is
    \[
    0\equiv +1\text{ eigenvalue},
    \qquad
    1\equiv -1\text{ eigenvalue}.
    \]
    
    For logical quantum views \(u,v\), write
    \[
    \operatorname{CNOT}_{u\rightarrow v}
    \]
    for the selected-qubit CNOT update with ordinary liveness, disjointness,
    and dimension obligations.
    
    If
    \[
    \operatorname{supp}(h)=\{t_1<\cdots<t_w\},
    \]
    then the \(Z(h)\)-parity measurement
    \[
    s:=\mathsf{MeasZCheck}_{h}(\bar r;a)
    \]
    is syntactic sugar for
    \[
    a:=\ket0;
    \quad
    \operatorname{CNOT}_{\bar r[t_1]\rightarrow a};
    \quad\cdots\quad;
    \operatorname{CNOT}_{\bar r[t_w]\rightarrow a};
    \quad
    s:=\mathcal M_{\mathrm{com}}[a].
    \]
    If \(w=0\), the CNOT sequence is empty and the measurement deterministically
    returns \(s=0\).
    
    The \(X(h)\)-parity measurement
    \[
    s:=\mathsf{MeasXCheck}_{h}(\bar r;a)
    \]
    is syntactic sugar for
    \[
    a:=\ket0;
    \quad
    a\asteq H;
    \quad
    \operatorname{CNOT}_{a\rightarrow \bar r[t_1]};
    \quad\cdots\quad;
    \operatorname{CNOT}_{a\rightarrow \bar r[t_w]};
    \quad
    a\asteq H;
    \quad
    s:=\mathcal M_{\mathrm{com}}[a].
    \]
    If \(w=0\), the CNOT sequence is empty and the measurement deterministically
    returns \(s=0\).
    
    These expand to initialization, Hadamard, CNOT, sequencing, and binary
    measurement before denotation.
    
    Let
    \[
    H_Z\in\{0,1\}^{r_Z\times N},
    \qquad
    H_X\in\{0,1\}^{r_X\times N}
    \]
    be fixed external CSS parity-check matrices satisfying
    \[
    H_XH_Z^{\mathsf T}=0
    \quad\text{over }\mathbb F_2.
    \]
    Write \(h^Z_u\) for the \(u\)-th row of \(H_Z\), and \(h^X_v\) for the
    \(v\)-th row of \(H_X\).
    
    For syndrome output registers
    \[
    s_Z[1,\ldots,r_Z],
    \qquad
    s_X[1,\ldots,r_X],
    \]
    and ancilla blocks
    \[
    a_Z[1,\ldots,r_Z],
    \qquad
    a_X[1,\ldots,r_X],
    \]
    define
    \[
    (s_Z,s_X):=
    \mathsf{CSSSyn}_{H_Z,H_X}(\bar r;a_Z,a_X)
    \]
    as syntactic sugar for
    \[
    \mathbf{for}^{r_Z}\ u=1\ \mathbf{to}\ r_Z
    \quad
    s_Z[u]:=\mathsf{MeasZCheck}_{h^Z_u}(\bar r;a_Z[u])
    \quad
    \mathbf{end};
    \]
    \[
    \mathbf{for}^{r_X}\ v=1\ \mathbf{to}\ r_X
    \quad
    s_X[v]:=\mathsf{MeasXCheck}_{h^X_v}(\bar r;a_X[v])
    \quad
    \mathbf{end}.
    \]
    
    Obligations: \(\bar r,a_Z,a_X\) are live quantum views; ancillas are
    pairwise disjoint and disjoint from \(\bar r\); \(s_Z,s_X\) are live bits.
    \(H_XH_Z^{\mathsf T}=0\) is needed for joint-syndrome projectors, not for
    parsing.

    \item \textbf{Finite for-loop notation.}

    For any fixed protocol length \(M\) and any statement schema \(S(i)\), we
    write
    \[
    \mathbf{for}^{M}\ i=1\ \mathbf{to}\ M
    \quad S(i)
    \quad \mathbf{end}
    \]
    as the finite expansion
    \[
    S(1);S(2);\cdots;S(M).
    \]
    This is fixed finite sequencing; it is not a while-loop.
    
    \item \textbf{CSS-QEC syntax sugar.}

    The abbreviation
    \[
    \bar u:=QEC(\bar r)
    \]
    abbreviates CSS correction sugar:
    syndrome measurement, decoder lookups, indexed Pauli corrections, and CSS
    decoding.

    \item \(\mathbf{rand}(n,\varphi)\) and \(\mathbf{rand}(n)\):
    
    The primitive random command is
    \[
    p:=_{\$}\mathbf{rand}(n,\varphi),
    \]
    where \(\varphi\) is a filter expression satisfying
    \[
    \llbracket\varphi\rrbracket_{\xi}\in\{0,1\}^{n}
    \]
    on well-formed branches.
    
    Its one-program denotation uniformly initializes the register payloads
    \[
    p[1],\ldots,p[n].
    \]
    The filter expression \(\varphi\) is descriptor data. It is used by relational
    rules to certify that the two samplings are aligned under the same local view
    index. It does not change the one-program randomization denotation.
    
    The notation
    \[
    p:=_{\$}\mathbf{rand}(n)
    \]
    is syntactic sugar for
    \[
    p:=_{\$}\mathbf{rand}(n,\mathbf 1_n),
    \]
    where \(\mathbf 1_n\) is the deterministic all-one filter of length \(n\).
    
    
    \item \(\mathbf{select}_{f}\):
    
    Let \(g\) be an indexed logical descriptor view, and let \(g_0,g_1\) be
    logical lists with
    \[
    \operatorname{sort}(g_0)=\operatorname{sort}(g_1)=\operatorname{sort}(g).
    \]
    Let \(f\) be a bit-valued filter descriptor.  The command
    \[
    (g_0,g_1)
    :=
    \mathbf{select}_{f}^{N}\ \mathbf{from}\ J\ \mathbf{for}\ g
    \]
    is well formed only when \(J\) is the whole current length of \(g\):
    \[
    J=\mathbf{length}(g)
    \]
    or \(J=N\) with
    \[
    N=\mathbf{length}(g).
    \]
    
    It defines \(g_0\) as the \(f=0\) subview of \(g\), and \(g_1\) as the
    \(f=1\) subview of \(g\), preserving relative order in each fiber.
    
    The package notation
    \[
    (\mathcal P_0,\mathcal P_1)
    :=
    \mathbf{select}_{f}^{N}\ \mathbf{from}\ J\ \mathbf{for}\ \mathcal P
    \]
    is componentwise.
    
    \(\mathbf{select}_{f}\) constructs descriptor views, not value/state copies.

    \item \(\mathbf{merge}_{f}\):
    
    Let \(g\) be a target descriptor view, let \(g_0,g_1\) be source views, and let
    \(f\) be a bit-valued filter descriptor.  The command
    \[
    g
    :=
    \mathbf{merge}_{f}^{N}\ \mathbf{from}\ J\ \mathbf{for}\ g_0,g_1
    \]
    is well formed only when \(J\) is the whole merged length:
    \[
    J=\mathbf{length}(g_0)+\mathbf{length}(g_1)
    \]
    or \(J=N\) with
    \[
    N=\mathbf{length}(g_0)+\mathbf{length}(g_1).
    \]
    
    It interleaves \(g_0,g_1\) according to \(f\), preserving source order.
    
    The package notation
    \[
    \mathcal P
    :=
    \mathbf{merge}_{f}^{N}\ \mathbf{from}\ J\ \mathbf{for}\ \mathcal P_0,\mathcal P_1
    \]
    is componentwise.
    
    \(\mathbf{merge}_{f}\) constructs descriptor views, not value/state copies.



    \item {\textbf{Bit-order send syntax sugar.}}
    
    {A channel attack is one fixed joint unitary}
    \[
    {
    U_{\textbf{channel}},
    }
    \]
    {shared by compared games. The only send sugar is}
    \[
    {
    \mathbf{send}_{\mathsf{bit}}(\vec u),
    \qquad
    \vec u=(u_1,\ldots,u_M),
    }
    \]
    {where \(u_1,\ldots,u_M\) is the displayed bit-order tuple.
    It expands to}
    \[
    {
    \mathbf{send}_{\mathsf{bit}}(u_1,\ldots,u_M)
    \stackrel{\triangle}{=}
    \bigl((u_1,\ldots,u_M),\bar q_{\mathbf{ch}}\bigr)
    \asteq
    U_{\textbf{channel}} .
    }
    \]
    
    {\(\bar q_{\mathbf{ch}}\) is the fixed channel/adversarial
    workspace. Obligations: \(u_1,\ldots,u_M\) are live, pairwise disjoint,
    dimension-compatible, and disjoint from that workspace.}
    
    {The displayed order is the expansion order.
    \(U_{\textbf{channel}}\) may entangle all transmitted systems with the
    workspace; it is not assumed product or decomposed.}
    
    {This sugar does not read/branch on filters, define a family
    of adversarial unitaries, or use dummy padding. If the tuple was selected,
    selection has already been represented by descriptor construction.}
    
    \item \textbf{publish:}

    For transcript \(L\) and classical/probability
    \(p=(p[1],\ldots,p[n])\), define
    \[
    \mathbf{publish: }p
    \]
    as value-copy publication:
    \[
    \textbf{for } i=1 \textbf{ to } i=n
    \quad
    L[\mathbf{next}] :=_{\mathsf{pub}} p[i]
    \quad
    \textbf{end}.
    \]
    We define the operation
    \[
    L[\mathbf{next}] :=_{\mathsf{pub}} p[i]
    \]
    to copy \(p[i]\)'s current value into the next reserved transcript cell,
    without descriptor aliasing or consuming the source.

    \(L\) is public/authenticated: visible but adversary-immutable.
    Publication is value-copy for classical/probability values only; quantum
    values cannot be value-copied. The source remains live.

    The displayed command
    \[
    \mathbf{publish:}~received
    \]
    is deterministic public acknowledgement after the channel command:
    \[
    received[1] := 1;\quad
    L[\mathbf{next}] :=_{\mathsf{pub}} received[1].
    \]
    It adds no loss model, availability oracle, or channel semantics;
    availability loss is represented by explicit later events such as
    \[
    \mathsf{sift\mbox{-}fail}.
    \]
\end{itemize}



\end{document}
